% The optimal CGLMP measurements on maximally entangled states.
% Physical Superintelligence (PSI), draft of 21 October 2026.
% Compile with: latexmk -pdf oqp27_prl_draft.tex
% (pdflatex + bibtex; the bibliography is embedded in the filecontents block).

\documentclass[aps,prl,reprint,superscriptaddress,longbibliography,floatfix,raggedbottom]{revtex4-2}
\usepackage[T1]{fontenc}
\usepackage{amsmath,amssymb,amsthm}
\usepackage{newtxtext,newtxmath}
\usepackage{microtype}
\usepackage{booktabs}
\usepackage{xcolor}
\usepackage[hidelinks]{hyperref}
\usepackage{filecontents}
\hypersetup{
  pdftitle={The optimal CGLMP measurements on maximally entangled states},
  pdfauthor={David Turturean and Ning Bao}
}
\allowdisplaybreaks[1]
\setcounter{secnumdepth}{3}

\newtheorem{theorem}{Theorem}
\newtheorem{lemma}{Lemma}
\newtheorem{corollary}{Corollary}

\newcommand{\C}{\mathbb C}
\newcommand{\R}{\mathbb R}
\newcommand{\Z}{\mathbb Z}
\newcommand{\T}{\mathbb T}
\newcommand{\ii}{\mathrm i}
\newcommand{\dd}{\,\mathrm d}
\newcommand{\ind}{\mathbf 1}
\newcommand{\ket}[1]{\lvert #1\rangle}
\newcommand{\bra}[1]{\langle #1\rvert}
\newcommand{\norm}[1]{\lVert #1\rVert}
\newcommand{\abs}[1]{\lvert #1\rvert}
\newcommand{\ip}[2]{\langle #1,#2\rangle}
\DeclareMathOperator{\Tr}{Tr}
\DeclareMathOperator{\rank}{rank}
\DeclareMathOperator{\spec}{spec}
\DeclareMathOperator{\diag}{diag}
\DeclareMathOperator{\Ran}{Ran}
\DeclareMathOperator{\sech}{sech}
\DeclareMathOperator{\Beta}{Beta}
\DeclareMathOperator*{\Res}{Res}
\DeclareMathOperator{\Ex}{\mathbb E}
\DeclareMathOperator{\Prob}{Pr}
\newcommand{\tr}{\tau}
\newcommand{\IF}{I_{\mathrm F}}
\newcommand{\Gcat}{G}
\newcommand{\cstar}{c_*}
\newcommand{\Fd}{\mathfrak F_d}
\newcommand{\Gd}{\mathcal G_d}
\newcommand{\Hilb}{\mathsf H}
\newcommand{\Zisom}{\mathcal Z}
\newcommand{\Comp}[1]{\mathsf C^{(#1)}}
\newcommand{\RootP}{\mathcal M}
\providecommand{\doi}[1]{\href{https://doi.org/#1}{\nolinkurl{doi:#1}}}
\providecommand{\arxiv}[1]{\href{https://arxiv.org/abs/#1}{\nolinkurl{arXiv:#1}}}
\makeatletter
\protected\def\bibfnamefont#1{\in@{-}{#1}\ifin@\mbox{#1}\else#1\fi}
\makeatother

\begin{filecontents*}[overwrite]{oqp27-note-references.bib}
@article{CGLMP,
 author = {D. Collins and N. Gisin and N. Linden and S. Massar and S. Popescu},
 title = {Bell inequalities for arbitrarily high-dimensional systems},
 journal = {Phys. Rev. Lett.}, volume = {88}, pages = {040404}, year = {2002},
 note = {\doi{10.1103/PhysRevLett.88.040404}; \arxiv{quant-ph/0106024}}}
@article{KGZMZ,
 author = {D. Kaszlikowski and P. Gnaci{\'n}ski and M. {\.Z}ukowski and W. Miklaszewski and A. Zeilinger},
 title = {Violations of local realism by two entangled {$N$}-dimensional systems are stronger than for two qubits},
 journal = {Phys. Rev. Lett.}, volume = {85}, pages = {4418--4421}, year = {2000},
 note = {\doi{10.1103/PhysRevLett.85.4418}; \arxiv{quant-ph/0005028}}}
@article{DKZ,
 author = {T. Durt and D. Kaszlikowski and M. {\.Z}ukowski},
 title = {Violations of local realism with quantum systems described by {$N$}-dimensional {Hilbert} spaces up to {$N=16$}},
 journal = {Phys. Rev. A}, volume = {64}, pages = {024101}, year = {2001},
 note = {\doi{10.1103/PhysRevA.64.024101}; \arxiv{quant-ph/0101084}}}
@misc{GillProblem,
 author = {R. D. Gill},
 title = {The power of {CGLMP} inequalities},
 howpublished = {Problem 27 of R.~F. Werner's list \emph{Open problems in quantum information theory}, TU Braunschweig, posted 15 April 2005; now maintained by IQOQI Vienna, \url{https://oqp.iqoqi.oeaw.ac.at/the-power-of-cglmp-inequalities}},
 year = {2005}}
@misc{KruegerWerner,
 author = {O. Krueger and R. F. Werner},
 title = {Some open problems in quantum information theory},
 howpublished = {arXiv preprint}, year = {2005},
 note = {\arxiv{quant-ph/0504166}}}
@article{BGP,
 author = {J.-D. Bancal and N. Gisin and S. Pironio},
 title = {Looking for symmetric {Bell} inequalities},
 journal = {J. Phys. A: Math. Theor.}, volume = {43}, pages = {385303}, year = {2010},
 note = {\doi{10.1088/1751-8113/43/38/385303}; \arxiv{1004.4146}}}
@article{Fine,
 author = {A. Fine},
 title = {Hidden variables, joint probability, and the {Bell} inequalities},
 journal = {Phys. Rev. Lett.}, volume = {48}, pages = {291--295}, year = {1982},
 note = {\doi{10.1103/PhysRevLett.48.291}}}
@article{VDGG,
 author = {W. van Dam and R. D. Gill and P. D. Gr{\"u}nwald},
 title = {The statistical strength of nonlocality proofs},
 journal = {IEEE Trans. Inform. Theory}, volume = {51}, pages = {2812--2835}, year = {2005},
 note = {\doi{10.1109/TIT.2005.851738}; \arxiv{quant-ph/0307125}}}
@article{AGG,
 author = {A. Ac{\'i}n and R. Gill and N. Gisin},
 title = {Optimal {Bell} tests do not require maximally entangled states},
 journal = {Phys. Rev. Lett.}, volume = {95}, pages = {210402}, year = {2005},
 note = {\doi{10.1103/PhysRevLett.95.210402}; \arxiv{quant-ph/0506225}}}
@misc{Mishra26a,
 author = {A. Mishra and A. Senthilkumar},
 title = {{CGLMP} on maximally entangled states: exact optimality, rigidity and {Tsirelson} bounds},
 howpublished = {Code, certificates and manuscript, version 1.0, 30 September 2026, \url{https://github.com/anshM123/CGLMP}},
 year = {2026}}
@misc{Mishra26b,
 author = {A. Mishra and A. Senthilkumar},
 title = {Optimal {CGLMP} measurements in every dimension and the clauses of {IQOQI} \mbox{Vienna} {Open} {Quantum} {Problem} {27B}},
 howpublished = {Preprint, code and certificates; version 1.0 released 2 October 2026. Comparison uses \href{https://github.com/anshM123/IQOQI-OQP-27/tree/d5213b7b1c9975b8cc0f16d55e3baa9c4ba451a0}{commit \texttt{d5213b7}}, 20 October 2026, \url{https://github.com/anshM123/IQOQI-OQP-27}},
 year = {2026}}
@misc{Zhang26,
 author = {Y. Zhang},
 title = {{CGLMP} measurements need not maximize statistical strength},
 howpublished = {Manuscript no.~15 in \emph{Agentic proofs for QIQC: collected manuscripts}, version 1.1, Zenodo, 28 September 2026},
 year = {2026},
 note = {\doi{10.5281/zenodo.23022433}}}
@article{Masanes,
 author = {{\relax Ll}. Masanes},
 title = {Tight {Bell} inequality for {$d$}-outcome measurements correlations},
 journal = {Quantum Inf. Comput.}, volume = {3}, pages = {345--358}, year = {2003},
 note = {\doi{10.26421/QIC3.4-4}; \arxiv{quant-ph/0210073}}}
@article{ADGL,
 author = {A. Ac{\'i}n and T. Durt and N. Gisin and J. I. Latorre},
 title = {Quantum nonlocality in two three-level systems},
 journal = {Phys. Rev. A}, volume = {65}, pages = {052325}, year = {2002},
 note = {\doi{10.1103/PhysRevA.65.052325}; \arxiv{quant-ph/0111143}}}
@article{ZG,
 author = {S. Zohren and R. D. Gill},
 title = {Maximal violation of the {Collins--Gisin--Linden--Massar--Popescu} inequality for infinite dimensional states},
 journal = {Phys. Rev. Lett.}, volume = {100}, pages = {120406}, year = {2008},
 note = {\doi{10.1103/PhysRevLett.100.120406}; \arxiv{quant-ph/0612020}}}
@article{ZRGW,
 author = {S. Zohren and P. Reska and R. D. Gill and W. Westra},
 title = {A tight {Tsirelson} inequality for infinitely many outcomes},
 journal = {EPL (Europhys. Lett.)}, volume = {90}, pages = {10002}, year = {2010},
 note = {\doi{10.1209/0295-5075/90/10002}; \arxiv{1003.0616}}}
@article{LVN,
 author = {B. Lang and T. V{\'e}rtesi and M. Navascu{\'e}s},
 title = {Closed sets of correlations: answers from the zoo},
 journal = {J. Phys. A: Math. Theor.}, volume = {47}, pages = {424029}, year = {2014},
 note = {\doi{10.1088/1751-8113/47/42/424029}; \arxiv{1402.2850}}}
@article{NPA,
 author = {M. Navascu{\'e}s and S. Pironio and A. Ac{\'i}n},
 title = {Bounding the set of quantum correlations},
 journal = {Phys. Rev. Lett.}, volume = {98}, pages = {010401}, year = {2007},
 note = {\doi{10.1103/PhysRevLett.98.010401}; \arxiv{quant-ph/0607119}}}
@misc{IR,
 author = {M. Ioannou and D. Rosset},
 title = {Noncommutative polynomial optimization under symmetry},
 howpublished = {arXiv preprint}, year = {2021},
 note = {\arxiv{2112.10803}}}
@article{Yang,
 author = {T. H. Yang and T. V{\'e}rtesi and J.-D. Bancal and V. Scarani and M. Navascu{\'e}s},
 title = {Robust and versatile black-box certification of quantum devices},
 journal = {Phys. Rev. Lett.}, volume = {113}, pages = {040401}, year = {2014},
 note = {\doi{10.1103/PhysRevLett.113.040401}; \arxiv{1406.7127}}}
@article{SATWAP,
 author = {A. Salavrakos and R. Augusiak and J. Tura and P. Wittek and A. Ac{\'i}n and S. Pironio},
 title = {Bell inequalities tailored to maximally entangled states},
 journal = {Phys. Rev. Lett.}, volume = {119}, pages = {040402}, year = {2017},
 note = {\doi{10.1103/PhysRevLett.119.040402}; \arxiv{1607.04578}}}
@article{Sarkar,
 author = {S. Sarkar and D. Saha and J. Kaniewski and R. Augusiak},
 title = {Self-testing quantum systems of arbitrary local dimension with minimal number of measurements},
 journal = {npj Quantum Inf.}, volume = {7}, pages = {151}, year = {2021},
 note = {\doi{10.1038/s41534-021-00490-3}; \arxiv{1909.12722}}}
@article{Kaniewski,
 author = {J. Kaniewski and I. {\v{S}}upi{\'c} and J. Tura and F. Baccari and A. Salavrakos and R. Augusiak},
 title = {Maximal nonlocality from maximal entanglement and mutually unbiased bases, and self-testing of two-qutrit quantum systems},
 journal = {Quantum}, volume = {3}, pages = {198}, year = {2019},
 note = {\doi{10.22331/q-2019-10-24-198}; \arxiv{1807.03332}}}
@article{Tavakoli21,
 author = {A. Tavakoli and M. Farkas and D. Rosset and J.-D. Bancal and J. Kaniewski},
 title = {Mutually unbiased bases and symmetric informationally complete measurements in {Bell} experiments},
 journal = {Sci. Adv.}, volume = {7}, pages = {eabc3847}, year = {2021},
 note = {\doi{10.1126/sciadv.abc3847}; \arxiv{1912.03225}}}
@article{SupicBowles,
 author = {I. {\v{S}}upi{\'c} and J. Bowles},
 title = {Self-testing of quantum systems: a review},
 journal = {Quantum}, volume = {4}, pages = {337}, year = {2020},
 note = {\doi{10.22331/q-2020-09-30-337}; \arxiv{1904.10042}}}
@article{Brunner,
 author = {N. Brunner and D. Cavalcanti and S. Pironio and V. Scarani and S. Wehner},
 title = {Bell nonlocality},
 journal = {Rev. Mod. Phys.}, volume = {86}, pages = {419--478}, year = {2014},
 note = {\doi{10.1103/RevModPhys.86.419}; \arxiv{1303.2849}}}
@article{Bell,
 author = {J. S. Bell},
 title = {On the {Einstein} {Podolsky} {Rosen} paradox},
 journal = {Physics Physique Fizika}, volume = {1}, pages = {195--200}, year = {1964},
 note = {\doi{10.1103/PhysicsPhysiqueFizika.1.195}}}
@article{CHSH,
 author = {J. F. Clauser and M. A. Horne and A. Shimony and R. A. Holt},
 title = {Proposed experiment to test local hidden-variable theories},
 journal = {Phys. Rev. Lett.}, volume = {23}, pages = {880--884}, year = {1969},
 note = {\doi{10.1103/PhysRevLett.23.880}}}
@article{Tsirelson,
 author = {B. S. Cirel'son},
 title = {Quantum generalizations of {Bell}'s inequality},
 journal = {Lett. Math. Phys.}, volume = {4}, pages = {93--100}, year = {1980},
 note = {\doi{10.1007/BF00417500}}}
@article{Vaziri,
 author = {A. Vaziri and G. Weihs and A. Zeilinger},
 title = {Experimental two-photon, three-dimensional entanglement for quantum communication},
 journal = {Phys. Rev. Lett.}, volume = {89}, pages = {240401}, year = {2002},
 note = {\doi{10.1103/PhysRevLett.89.240401}}}
@article{Dada,
 author = {A. C. Dada and J. Leach and G. S. Buller and M. J. Padgett and E. Andersson},
 title = {Experimental high-dimensional two-photon entanglement and violations of generalized {Bell} inequalities},
 journal = {Nature Phys.}, volume = {7}, pages = {677--680}, year = {2011},
 note = {\doi{10.1038/nphys1996}; \arxiv{1104.5087}}}
@article{Miao26,
 author = {J.-L. Miao and E. Svegborn and Z. Chen and Y. Guo and X.-M. Hu and Y.-F. Huang and C.-F. Li and G.-C. Guo and A. Tavakoli and B.-H. Liu},
 title = {Binarization-loophole-free observation of high-dimensional quantum nonlocality},
 journal = {Phys. Rev. Lett.}, volume = {136}, pages = {240201}, year = {2026},
 note = {\doi{10.1103/ydch-vgmy}; \arxiv{2601.02350}}}
@book{Katznelson,
 author = {Y. Katznelson},
 title = {An Introduction to Harmonic Analysis},
 edition = {3rd}, publisher = {Cambridge University Press}, address = {Cambridge}, year = {2004},
 note = {\doi{10.1017/CBO9781139165372}}}
@misc{Heinavaara26,
 author = {O. Hein{\"a}vaara},
 title = {Convolution comparison measures},
 howpublished = {arXiv preprint}, year = {2026},
 note = {\arxiv{2602.10373}}}
@misc{DLMF,
 author = {{F. W. J. Olver and others (eds.)}},
 title = {{NIST Digital Library of Mathematical Functions}, {\S25.11}},
 howpublished = {\url{https://dlmf.nist.gov/25.11}, release 1.2.8 of 15 September 2026},
 year = {2026}}
@article{Bernstein,
 author = {S. N. Bernstein},
 title = {D{\'e}monstration du th{\'e}or{\`e}me de {Weierstrass} fond{\'e}e sur le calcul des probabilit{\'e}s},
 journal = {Comm. Soc. Math. Kharkov (2)}, volume = {13}, pages = {1--2}, year = {1912}}
@article{GoldsteinReinert,
 author = {L. Goldstein and G. Reinert},
 title = {Stein's method for the {Beta} distribution and the {P{\'o}lya--Eggenberger} urn},
 journal = {J. Appl. Probab.}, volume = {50}, pages = {1187--1205}, year = {2013},
 note = {\doi{10.1239/jap/1389370107}; \arxiv{1207.1460}}}
@article{Dobler,
 author = {C. D{\"o}bler},
 title = {Stein's method of exchangeable pairs for the {Beta} distribution and generalizations},
 journal = {Electron. J. Probab.}, volume = {20}, number = {109}, pages = {1--34}, year = {2015},
 note = {\doi{10.1214/EJP.v20-3933}; \arxiv{1411.4477}}}
@book{Moore,
 author = {R. E. Moore and R. B. Kearfott and M. J. Cloud},
 title = {Introduction to Interval Analysis},
 publisher = {SIAM}, address = {Philadelphia}, year = {2009},
 note = {\doi{10.1137/1.9780898717716}}}
\end{filecontents*}

\AtBeginDocument{\providecommand{\doi}{}\renewcommand{\doi}[1]{\href{https://doi.org/#1}{\nolinkurl{doi:#1}}}}
\begin{document}

\title{The optimal CGLMP measurements on maximally entangled states}
\author{David Turturean}\thanks{D.T. built and ran the automated research and formalization workflows at PSI and steered the writing and revision of the paper.}
\affiliation{Physical Superintelligence PBC, Cambridge, Massachusetts, USA}
\author{Ning Bao}\thanks{N.B. commented on the manuscript, as a paid contractor to PSI, not under one of his grants or as an employee of Northeastern University.}
\affiliation{Northeastern University}
\affiliation{Brookhaven National Laboratory}

\begin{abstract}
For projective measurements on a prescribed maximally entangled state,
we prove the optimal CGLMP value and classify every maximizer. Equality
requires the local dimension to be a multiple of the number of measurement outcomes, and
the measurements are Fourier-phase bases.
A positive pinching identity and a tilted random partition give a deficit
certificate.
This resolves the first assertion of Problem 27.B in the projective class.
\end{abstract}

\maketitle

{\raggedright\setlength{\parskip}{0.4\baselineskip}
\section*{Acknowledgements}

We thank Adam Levine for revising the manuscript, Ian MacCormack for reviewing earlier drafts and circulating a concise version of the manuscript, and Jim Halverson for coordinating the review of the manuscript. We also thank Physical Superintelligence PBC, in which some of the authors have a financial interest, for use of its resources, including its technology and frontier AI models. The authors used these models to develop, derive, and compose this manuscript and its Lean~4 supplements.

Results were formalized in Lean~4 by AI agents and checked by Lean's kernel; the formalization is at \url{https://github.com/physical-superintelligence/oqp-batch-2026/tree/main/oqp27-cglmp-inequalities/lean}.

David Turturean and Ning Bao reviewed and edited this paper and are responsible for its correctness. David Turturean is employed by PSI and has a financial interest in it.

\par}
\par\medskip\noindent\rule{\columnwidth}{0.4pt}\par\medskip


Bell inequalities distinguish quantum correlations from those produced
by local hidden variables \cite{Bell,CHSH,Brunner}. The CGLMP family
\cite{CGLMP} uses two measurement settings per party and $d$ outcomes.
For the prescribed maximally entangled state, which projective
measurements give the largest violation? Numerical searches
\cite{KGZMZ,DKZ} found Fourier-phase measurements in every dimension
tested. We prove their optimality and classify all maximizers, allowing
arbitrary projection ranks and arbitrary local dimension. This resolves
the first assertion of Problem~27.B for projective measurements on
maximally entangled states \cite{GillProblem,KruegerWerner}.
Section~\ref{sec:statement} states the result;
Sections~\ref{sec:reduction}--\ref{sec:equality} prove it.

\section{\texorpdfstring{The CGLMP expression and the theorem}{The CGLMP expression and the theorem}}\label{sec:statement}

\emph{Bell scenarios and the local bound.---}Alice and Bob choose
settings $x,y\in\{0,1\}$ and obtain outcomes $a,b\in\Z_d$. Their
behaviour is the conditional distribution $q(a,b\,|\,x,y)$. Local
behaviours are convex combinations of deterministic assignments
$(a_0,a_1,b_0,b_1)\in\Z_d^4$; these form the local polytope.
For $d=2$, recording outcomes as signs gives the CHSH inequality~\cite{CHSH},
\begin{equation}\label{eq:chsh}
\langle A_0B_0\rangle+\langle A_1B_0\rangle
+\langle A_1B_1\rangle-\langle A_0B_1\rangle\le2.
\end{equation}
Its $d$-outcome generalization uses the elementary identity
\begin{equation}\label{eq:parity}
(a_0-b_0)+(b_0-a_1)+(a_1-b_1)+(b_1-a_0-1)=-1.
\end{equation}

\emph{The CGLMP inequality.---}For outcomes in $\Z_d$ define the residuals
\[
\begin{gathered}
r_{00}=[a-b]_d,\qquad r_{10}=[b-a]_d,\\
r_{11}=[a-b]_d,\qquad r_{01}=[b-a-1]_d,
\end{gathered}
\]
where $[u]_d\in\{0,\dots,d-1\}$ is $u$ mod $d$,
and put
\begin{equation}\label{eq:Id}
S_d(q)=\sum_{x,y\in\{0,1\}}\Ex_q\bigl(r_{xy}\bigr),\qquad
I_d(q)=4-\frac{2S_d(q)}{d-1}.
\end{equation}
For a deterministic assignment the sum of the four residuals is congruent,
by \eqref{eq:parity}, to $-1$ modulo $d$, so this sum---a sum of four
representatives in $\{0,\dots,d-1\}$---is at least $d-1$. Hence
$S_d\ge d-1$, that is $I_d\le2$, for every local behaviour. This is
the CGLMP inequality \cite{CGLMP}, in the normalization of \cite{AGG},
chosen so that the local bound is $2$ in every dimension; at $d=2$,
unwinding the residuals shows that $I_2$ is exactly the left side of
\eqref{eq:chsh}. The CGLMP inequality is a facet of the local polytope for every $d$
\cite{Masanes}. Its value alone does not determine statistical strength
or resistance to an unspecified noise model; we distinguish those
quantities in the Discussion.

\emph{Quantum strategies.---}A \emph{strategy} is any physical
arrangement that produces a behaviour; the strategies we study are
quantum. Alice and Bob share the maximally entangled state
\[
\Phi_n=\frac1{\sqrt n}\sum_{j=0}^{n-1}\ket{jj}
\]
of two $n$-level systems, and each measurement is projective: it is
described by a \emph{projection-valued measure} (PVM), a family of $d$
orthogonal projections on $\C^n$ summing to the identity, one per
outcome, the probability of outcome $a$ being the expectation of the
$a$th projection. Some projections may be zero, and their ranks may
differ. A \emph{projective strategy} on $\Phi_n$ is a choice of four
such PVMs, two for Alice and two for Bob. Writing $A^x_a$
for the projection of Alice's measurement $x$ onto outcome $a$, and
$B^y_b$ likewise for Bob, the behaviour is
\[
q(a,b\,|\,x,y)=\bra{\Phi_n}A^x_a\otimes B^y_b\ket{\Phi_n}
=\tr_n\bigl((A^x_a)^{T}B^y_b\bigr),
\]
where $\tr_n=\Tr/n$ is the normalized trace on $n\times n$ matrices
and the transpose is taken in the basis $\{\ket j\}$: on the maximally
entangled state, probabilities are normalized traces of products of
measurement operators. It is convenient to arrange the four PVMs in
the cyclic order
\begin{equation}\label{eq:tracial}
(P^0,P^1,P^2,P^3)=(A_0^{T},B_0,A_1^{T},B_1),
\end{equation}
and to write $P^t_a$ for the projection of the PVM $P^t$ onto outcome
$a\in\Z_d$, so that $q(a,b\,|\,x,y)=\tr_n(P^{2x}_aP^{1+2y}_b)$, with
the four measurement contexts corresponding to cyclically adjacent
pairs. Outcome labels are read modulo $d$ throughout. We call
$P^0,\dots,P^3$ the \emph{tracial PVMs} of the strategy.
Conjugating all four by one unitary $U$ on $\C^n$ corresponds to the
local unitary $\overline U\otimes U$, which leaves $\Phi_n$ unchanged;
strategies related in this way are physically equivalent.

\emph{The Fourier-phase measurements.---}For $t\in\{0,1,2,3\}$ and
$a\in\Z_d$ put
\begin{equation}\label{eq:standard}
\psi_{t,a}=\frac1{\sqrt d}\sum_{j=0}^{d-1}
e^{-2\pi\ii j(a-t/4)/d}\ket j,\quad
\Pi^t_a=\ket{\psi_{t,a}}\bra{\psi_{t,a}}.
\end{equation}
For each $t$ the vectors $\psi_{t,0},\dots,\psi_{t,d-1}$ form an
orthonormal basis of $\C^d$: the computational basis transformed by the
discrete Fourier transform and then by the diagonal unitary
$\diag(e^{\pi\ii jt/(2d)})_j$. Hence $\Pi^t=(\Pi^t_a)_{a\in\Z_d}$ is a
PVM for each $t$, and these four PVMs are the \emph{Fourier-phase
measurements}. The strategy with tracial PVMs
$P^t_a=\Pi^t_a$ on $\Phi_d$, so that Alice's two physical bases are the
complex conjugates of $\psi_{0,\cdot}$ and $\psi_{2,\cdot}$ and Bob's are
$\psi_{1,\cdot}$ and $\psi_{3,\cdot}$, has the CGLMP value
\begin{equation}\label{eq:IF}
I_d(q_{\mathrm F})  = \frac4{d(d-1)}\sum_{k=1}^{d-1}(d-k)\sec\frac{\pi k}{2d} \equiv \IF(d),
\end{equation}
where $q_{\mathrm F}$ is the behaviour of the Fourier-phase measurements on
$\Phi_d$.
The formula in \eqref{eq:IF} is Lemma~\ref{lem:companion} below with $V_j=I$. Thus
$\IF(2)=2\sqrt2$, the bound of Tsirelson \cite{Tsirelson},
$\IF(3)=2.8729\ldots$, and $\IF(d)\to32\Gcat/\pi^2=2.9698\ldots$ as
$d\to\infty$, where $\Gcat=\sum_{k\ge0}(-1)^k(2k+1)^{-2}=0.9159\ldots$ is
Catalan's constant \cite{DLMF,ZG}. These measurements were found
numerically by Kaszlikowski et al.\ \cite{KGZMZ} and by Durt,
Kaszlikowski and \.Zukowski \cite{DKZ}; the value \eqref{eq:IF} is
computed in \cite{CGLMP}. Our result is that no other choice of projective
measurements on a maximally entangled state gives a larger CGLMP
value, and that any choice attaining the same value is the
Fourier-phase measurements up to a local change of basis. Stated precisely we have the main theorem:

\begin{theorem}[Value and rigidity]\label{thm:main}
Let $d\ge2$. Every projective strategy on $\Phi_d$ satisfies
$I_d\le\IF(d)$. Equality holds if and only if there is a unitary $U$
on $\C^d$ such that the tracial PVMs are
\[
P^t_a=U\,\Pi^t_a\,U^\dagger\qquad\text{for all }t,a.
\]
\end{theorem}

In words: the maximizers are exactly the Fourier-phase measurements,
up to the local unitary $\overline U\otimes U$ acting on the maximally entangled state. In particular every
maximizing measurement is a measurement in a rank-one basis, although
nothing of the kind was assumed. Theorem~\ref{thm:general} in
Appendix~\ref{app:general} extends this to a maximally entangled state
of any local dimension $n$ but where the measurement outcomes still take one of $d$ values. Saturation of the inequality then requires $d$ to divide $n$ and
the maximizers are Fourier-phase measurements on a $d$-level subsystem
tensored with the identity of an inert $(n/d)$-level system. Furthermore, for
$d\nmid n$ the value $\IF(d)$ is not attained by any projective
strategy.

\emph{Previous work.---}The value \eqref{eq:IF} was conjectured to be
optimal on maximally entangled states since \cite{DKZ}. Lang, V\'ertesi
and Navascu\'es \cite{LVN} proved it numerically for $d=3$ to within
$10^{-10}$ with a semidefinite relaxation of the tracial problem
\eqref{eq:tracial}. Optimization over all states is a different problem:
partially entangled states violate the inequality more \cite{ADGL}, the
hierarchy of \cite{NPA} and the symmetry-adapted certificates of
\cite{IR} give the exact all-state value for $d=3,4$, Yang et al.\
\cite{Yang} certified the qutrit optimum in a black-box way, and
\cite{ZG,ZRGW} treated the limit $d\to\infty$. Device-independent
rigidity of the same measurements is known for Bell expressions tailored
to maximally entangled states: they are self-tested in every dimension
with two settings in \cite{SATWAP,Sarkar}, and mutually unbiased bases
are self-tested in \cite{Kaniewski,Tavakoli21}; see the reviews
\cite{Brunner,SupicBowles}. The Fourier-phase measurements are also the
ones used in high-dimensional Bell experiments \cite{Vaziri,Dada,Miao26}.
Mishra and Senthilkumar \cite{Mishra26b} prove the same projective value,
rigidity, and divisibility statements. The comparison here uses their
20 October 2026 commit, identified in the reference. Their clock
reduction uses a cyclic twirl and a finite Stone--von Neumann theorem;
Lemma~\ref{lem:companion} gives explicit transition matrices and the
reducing isometry. Both approaches use a positive matrix measure and a
strip inequality. Here Lemma~\ref{lem:pinching} gives a pinching identity
and its holomorphic continuation, using the trace-density moments of
Hein\"avaara \cite[Eq.~(16), proof of Lemma~2.5]{Heinavaara26}.
Random cells also occur in their argument; our partition law is the
explicit density \eqref{eq:partitiondensity}, bounded below by $c_*$.
Their cone certificate uses interval computations for $d\le200$,
$201\le d\le2000$, and a uniform computer-assisted estimate beyond
that range. Our beta--Stein estimate treats every $d\ge66$
analytically, leaving 2,080 scalar comparisons for $2\le d\le65$.
The contribution developed here is this explicit deficit certificate
and its direct forbidden-overlap route to equality classification.

\emph{Proof sketch.---}Section~\ref{sec:reduction} encodes every
strategy into $d$ unitary matrices $V_0,\dots,V_{d-1}$ with $V_j^4=I$,
in such a way that the deficit $\IF(d)-I_d$ becomes an explicit linear
expression in their pairwise overlaps $\tr(V_j^\dagger V_k)$. The task
is then to show that this expression is nonnegative for every such
family of matrices, and zero for essentially one. In
Section~\ref{sec:partition}, we use Fourier analysis to derive a
family of constraints on these overlaps, and combine them to show
that the deficit is bounded below by a
positive multiple of a manifestly nonnegative quantity, which measures
how far the eigenspaces of the different $V_j$ are from being aligned.
This proves the bound. In Section~\ref{sec:equality}, we
show that when the deficit vanishes, the $V_j$ must commute and take a single form that corresponds exactly to the
Fourier-phase measurements.

\emph{Conventions.---}$\tr=\tr_r=\Tr/r$ denotes the normalized trace on
$r\times r$ matrices; the size $r$ will be clear from context. For
matrices, $\ip XY=\tr(X^\dagger Y)$ and $\norm X_2^2=\tr(X^\dagger X)$;
$\norm X$ is the operator norm. A unitary $V$ with $V^4=I$ has the
spectral decomposition $V=\sum_{\ell\in\Z_4}\ii^{\ell}E_\ell$ with
orthogonal projections $E_0,\dots,E_3$ summing to $I$ (some possibly
zero); we call them the spectral projections of $V$ and read their index
modulo four.

\section{Encoding quantum strategies into unitaries}
\label{sec:reduction}

A projective strategy is specified by its four tracial PVMs
$P^0,P^1,P^2,P^3$ of \eqref{eq:tracial}, and only cyclically adjacent
pairs enter the CGLMP expression: the four contexts $(x,y)$ are the
pairs $(P^0,P^1)$, $(P^1,P^2)$, $(P^2,P^3)$ and $(P^3,P^0)$. It helps
to picture the four PVMs as the vertices of a square and the four
contexts as its edges. Lemma~\ref{lem:companion} below trades the four
PVMs for a single unitary of order four on an enlarged space---a
``clock'' that steps around the square---and then block-diagonalizes
that unitary into $d$ unitaries $V_0,\dots,V_{d-1}$ of order four, one
for each Fourier mode $j$ of the outcome label. All information about
$I_d$ survives in the normalized overlaps
\[
z_{jk}=\tr_r(V_j^\dagger V_k),\qquad 0\le j<k<d .
\]
For a coefficient sequence $c=(c_k)_{k=1}^{d-1}$ define the quadratic
form
\begin{equation}\label{eq:Lc}
L_c(V)=\sum_{j<k}
\bigl[c_{k-j}(1-\operatorname{Re}z_{jk})+c_{d-k+j}
\operatorname{Im}z_{jk}\bigr],
\end{equation}
and let
\begin{equation}\label{eq:Gdef}
\Gd(V)=L_\beta(V),\qquad \beta_k=\frac1{2d^2}\sec\frac{\pi k}{2d},
\end{equation}
so that explicitly
\[
\Gd=\frac1{2d^2}\sum_{j<k}\Bigl[\sec\tfrac{\pi(k-j)}{2d}
(1-\operatorname{Re}z_{jk})
+\csc\tfrac{\pi(k-j)}{2d}\operatorname{Im}z_{jk}\Bigr].
\]
The \emph{forbidden overlap} of the tuple is
\begin{equation}\label{eq:forbidden}
\Fd(V)=\sum_{j<k}p_{jk},\qquad
p_{jk}=\frac{1-\operatorname{Re}z_{jk}+\operatorname{Im}z_{jk}}2 .
\end{equation}
In terms of the spectral projections $V_j=\sum_\ell\ii^\ell E_{j\ell}$,
a short computation with
$z_{jk}=\sum_{\ell,m}\ii^{m-\ell}\tr(E_{j\ell}E_{km})$ gives
\begin{equation}\label{eq:forbiddenspectral}
p_{jk}=\sum_{\ell=0}^3\sum_{u=1}^2\tr(E_{j\ell}E_{k,\ell+u})\ge0,
\end{equation}
with spectral indices read modulo four: $p_{jk}$ collects the overlaps
in which the spectral label jumps by one or two units in the wrong
direction. Each term is nonnegative because
$\tr(EF)=\tr\bigl((EF)(EF)^\dagger\bigr)\ge0$ for projections $E,F$.
Thus $\Fd\ge0$, and $\Fd=0$ forces every forbidden product
$E_{j\ell}E_{k,\ell+1}$ and $E_{j\ell}E_{k,\ell+2}$ to vanish, which will be used to prove uniqueness in Section~\ref{sec:equality}.

The following construction retains the original PVMs on an explicit
reducing subspace, which is essential for the equality classification.

\begin{lemma}[Companion representation]\label{lem:companion}
Every projective strategy on $\Phi_n$, $n\ge1$, has a companion
tuple $V_0,\dots,V_{d-1}\in U(4n)$ with $V_j^4=I$ satisfying
\begin{equation}\label{eq:normalization}
\IF(d)-I_d=\frac{8d}{d-1}\,\Gd(V).
\end{equation}
Moreover, the tuple determines four PVMs
$\RootP^0(V),\dots,\RootP^3(V)$ on $\C^d\otimes\C^{4n}$, the
\emph{companion PVMs} \eqref{eq:rootmeasurements}, and an isometry
$\Zisom:\C^n\to\C^d\otimes\C^{4n}$ such that
\begin{equation}\label{eq:intertwine}
\RootP^t_a(V)\,\Zisom=\Zisom\,P^t_a\qquad\text{for all }t,a. 
\end{equation}
In other words, on the range of $\Zisom$, the companion PVMs act exactly as the
original ones.
\end{lemma}

\begin{proof}
\emph{Step 1: a random shift of the outcomes.}
Let $J\ket a=\ket{a-1}$ be the cyclic shift on $\C^d$, and define on
$\C^d\otimes\C^n$ the operators
\[
Z_t=\sum_{c=0}^{d-1}J^c\otimes P^t_c\quad(0\le t<4),
\qquad Z_4=Z_0(J\otimes I_n),
\]
where, as above, $P^t_c$ is the projection of $P^t$ onto outcome $c$.
Each $Z_t$ is unitary, because the $P^t_c$ are orthogonal projections
summing to $I_n$, and commutes with $J\otimes I_n$. Let
$\Comp n_a=\ket a\bra a\otimes I_n$ be the computational-basis PVM on
$\C^d\otimes\C^n$. Conjugating it by $Z_t$ gives
\[
Z_t\,\Comp n_a\,Z_t^\dagger
=\sum_{k=0}^{d-1}\ket k\bra k\otimes P^t_{a-k},
\]
with $a-k$ read modulo $d$. The four PVMs
$(Z_t\Comp n_aZ_t^\dagger)_{a\in\Z_d}$, $t=0,\dots,3$, form a
projective strategy on $\Phi_{dn}$ whose outcomes are those of the
original strategy shifted by a common, uniformly random $k\in\Z_d$:
for $t\le2$,
\[
\tr_{dn}\bigl(Z_t\Comp n_aZ_t^\dagger\,Z_{t+1}\Comp n_bZ_{t+1}^\dagger\bigr)
=\frac1d\sum_{k\in\Z_d}\tr_n\bigl(P^t_{a-k}P^{t+1}_{b-k}\bigr),
\]
and similarly on the last edge, where the extra shift in $Z_4$
accounts for the $-1$ in the residual $r_{01}$. The common shift
cancels from every residual $r_{xy}$, so $I_d$ is unchanged, and the
original strategy sits in the $k=0$ block.

\emph{Step 2: the clock.}
Put $\mathsf U_t=Z_t^\dagger Z_{t+1}$, the transition unitary along edge
$t$ of the square, from $P^t$ to $P^{t+1}$; the shift in $Z_4$
implements the $-1$ in the residual $r_{01}$. On a four-dimensional clock register define
\[
\mathcal R=\sum_{t=0}^3\ket t\bra{t+1\bmod4}\otimes\mathsf U_t .
\]
Every cyclic product $\mathsf U_t\mathsf U_{t+1}\mathsf U_{t+2}
\mathsf U_{t+3}$ telescopes to $Z_t^\dagger Z_t(J\otimes I_n)$, so
$\mathcal R^4=I_4\otimes J\otimes I_n$. Diagonalize $J$ by the discrete
Fourier transform $\mathsf F_{aj}=d^{-1/2}e^{2\pi\ii aj/d}$ and pick the
fourth roots $\lambda_j=e^{\pi\ii j/(2d)}$ of its eigenvalues. Since
$\mathcal R$ commutes with $I_4\otimes J\otimes I_n$, reordering the
tensor factors as $\C^d\otimes(\C^4\otimes\C^n)$ and Fourier
transforming the outcome coordinate gives
\begin{equation}\label{eq:blockdiag}
\mathcal R=(\mathsf F\otimes I_{4n})\,
\diag(\lambda_jV_j)_{j=0}^{d-1}\,
(\mathsf F^\dagger\otimes I_{4n})
\end{equation}
with $V_j\in U(4n)$; taking fourth powers shows $V_j^4=I$. In these
coordinates define the companion PVMs on $\C^d\otimes\C^{4n}$,
\begin{equation}\label{eq:rootmeasurements}
\RootP^t_a(V)=\mathcal R^t\,\Comp{4n}_a\,\mathcal R^{-t},
\qquad 0\le t<4 .
\end{equation}
Their value of $I_d$ on $\Phi_{4dn}$ equals the original: the edge
distribution
$\tr(\RootP^t_a\RootP^{t+1}_b)=\tr(\Comp{4n}_a\mathcal R\,
\Comp{4n}_b\mathcal R^\dagger)$ is independent of $t$, and tracing out
the clock averages the four transition unitaries $\mathsf U_t$, so each
companion edge carries one quarter of the four original edge costs.

The original strategy survives inside: in clock position zero,
$\mathsf U_0\cdots\mathsf U_{t-1}=Z_0^\dagger Z_t$, so the isometry
\begin{equation}\label{eq:reducingisometry}
\Zisom\xi=\sum_a\ket a\otimes\ket0\otimes P^0_a\xi
\end{equation}
satisfies $\Zisom^\dagger\Zisom=I_n$ and
$\RootP^t_a(V)\Zisom=\Zisom P^t_a$, which is \eqref{eq:intertwine}.
Since the $\RootP^t_a$ are self-adjoint, both the range of $\Zisom$ and
its orthogonal complement are invariant under every $\RootP^t_a$, and
on the range the companion PVMs are a copy of the original ones.

\emph{Step 3: the cost in terms of overlaps.}
Write $\omega=e^{2\pi\ii/d}$. By \eqref{eq:blockdiag}, the block of
$\mathcal R$ between computational levels $a$ and $b$ of the outcome
coordinate is $R_{ab}=d^{-1}\sum_j\omega^{j(a-b)}\lambda_jV_j$, which
depends only on $a-b$. The probability that the outcomes of one edge
differ by $c$ is therefore
\[
\Prob(c)=\frac1{d^2}\sum_{j,k=0}^{d-1}\omega^{c(j-k)}
\lambda_j\bar\lambda_k\,\tr_{4n}(V_k^\dagger V_j).
\]
Summing $\sum_c c\,\omega^{c(j-k)}=-d/(1-\omega^{j-k})$ for
$j\not\equiv k\pmod d$ and $\sum_c c=d(d-1)/2$ for $j=k$, and using,
for $x=\mbox{$(j-k)/d\in(0,1)$}$,
\begin{equation*}
\begin{split}
\frac{\lambda_j\bar\lambda_k}{1-\omega^{\,j-k}}
&=\frac{e^{\pi\ii x/2}}2\,(1+\ii\cot\pi x)\\
&=\frac14\Bigl(\sec\frac{\pi x}2+\ii\csc\frac{\pi x}2\Bigr),
\end{split}
\end{equation*}
the expected residual of one edge becomes
\begin{equation}\label{eq:edgecost}
\begin{split}
\frac{d-1}2-\frac1{2d}\sum_{j<k}\Bigl[
&\sec\tfrac{\pi(k-j)}{2d}\operatorname{Re}z_{jk}\\
&-\csc\tfrac{\pi(k-j)}{2d}\operatorname{Im}z_{jk}\Bigr].
\end{split}
\end{equation}
Subtracting the value of \eqref{eq:edgecost} at $V_j=I$ gives exactly
$d\,\Gd(V)$. Summing the four edges, $S_d$ exceeds its value at $V_j=I$
by $4d\,\Gd$, and \eqref{eq:Id} turns this into \eqref{eq:normalization},
provided the strategy at $V_j=I$ has value $\IF(d)$.
Indeed, put $R_0=\mathsf F\diag(\lambda_j)\mathsf F^\dagger$ and
$C_a=\ket a\bra a$. At $V_j=I$, the companion unitary is
$\mathcal R_0=R_0\otimes I_{4n}$. Since
$\mathsf F^\dagger\ket a=\psi_{0,a}$ and
$\diag(\lambda_j^t)\psi_{0,a}=\psi_{t,a}$, one has
\begin{equation}\label{eq:fourierconjugation}
R_0^t C_a R_0^{-t}=\mathsf F\Pi^t_a\mathsf F^\dagger.
\end{equation}
Thus the companion PVMs are a common Fourier conjugate of the standard
amplification. Their probabilities agree, and evaluating
\eqref{eq:edgecost} at $z_{jk}=1$ and inserting into \eqref{eq:Id}
yields \eqref{eq:IF}.
\end{proof}

\section{Bounding \texorpdfstring{$\boldsymbol{\Gd(V)}$}{G\_d(V)}}
\label{sec:partition}

Section~\ref{sec:reduction} reduced Theorem~\ref{thm:main} to a
statement about $d$ unitaries $V_0,\dots,V_{d-1}$ of order four: the
quantity $\Gd(V)$, a combination of their overlaps $z_{jk}$ with secant
and cosecant coefficients, must be nonnegative. We obtain inequalities
between the overlaps from the following construction. Cut the unit
circle, parametrized by $t\in[0,1)$, into $d$ consecutive arcs, put
$V(t)=V_j$ on the $j$th arc, and form the matrix-valued function
\[
U(e^{2\pi\ii t})=e^{-\pi\ii t/2}\,V(t)^\dagger .
\]
Because $V_j^4=I$, this function satisfies $U(z)^4=z^{-1}I$ at every
point $z=e^{2\pi\ii t}$ of the circle. (The phase $e^{-\pi\ii t/2}$ is
the continuous version of the fourth roots $\lambda_j=e^{\pi\ii j/(2d)}$
in \eqref{eq:blockdiag}: at $t=j/d$, $U$ equals $(\lambda_jV_j)^\dagger$.)
We study $U$ through its Fourier series.

Write $U(z)=\sum_{k\in\Z}c_kz^k$ in $L^2$, with
\[
c_k=\int_0^1U(e^{2\pi\ii t})e^{-2\pi\ii kt}\dd t .
\]
For the negative-frequency projection $P_-U=\sum_{k<0}c_kz^k$,
Parseval's identity gives
\[
\norm{P_-U}_2^2=\sum_{k<0}\tr_r(c_k^\dagger c_k).
\]
A unitary-valued $U$ has $\norm U_2^2=1$.

We are interested in unitary-valued $U$ with $U(z)^4=z^{-1}I$, which we
call \emph{fourth roots of $z^{-1}$}. Discontinuous ones are allowed
and necessary, since the functions built above from the $V_j$ have jumps at the partition
boundaries. The basic example has $r=1$:
$U_0(e^{2\pi\ii t})=e^{-\pi\ii t/2}$ for $0\le t<1$, which jumps by a
factor $-\ii$ when $t$ passes $1$. Its coefficients satisfy
$\abs{c_k}^2=8/[\pi^2(4k+1)^2]$. Using
$\sum_{n\ge0}(4n+3)^{-2}=\pi^2/16-\Gcat/2$, where
$\Gcat=\sum_{k\ge0}(-1)^k(2k+1)^{-2}=0.9159\ldots$ is Catalan's
constant from \eqref{eq:IF}, we find
\begin{equation}\label{eq:nu4}
\norm{P_-U_0}_2^2
=\frac8{\pi^2}\sum_{n\ge0}\frac1{(4n+3)^2}
=\frac12-\frac{4\Gcat}{\pi^2}=:\nu_4 .
\end{equation}
The next lemma says that no fourth root of $z^{-1}$, of any matrix
size, has a smaller value of $\norm{P_-U}_2^2$.

\begin{lemma}[Negative frequencies of fourth roots]\label{lem:hardy}
Let $U$ be a measurable function from the unit circle to the
$r\times r$ unitary matrices with $U(z)^4=z^{-1}I$ for almost every
$z$. Then
\begin{equation}\label{eq:Hardy}
\norm{P_-U}_2^2\ge\nu_4 .
\end{equation}
\end{lemma}


The proof is given in Appendix~\ref{app:hardy}, together with an
explanation of the idea. In brief, $1-2\norm{P_-U}_2^2$ is rewritten
as the average over the circle of $\tr_r(A\,\Hilb B)$, where $A$ and
$B$ are commuting Hermitian matrices with eigenvalues $\pm1$ built
from $U$, and $\Hilb B$ is obtained from $B$ by multiplying its $k$th
Fourier coefficient by $-\ii\operatorname{sgn}k$. The average is then
bounded using an analytic function $h$ on the strip
$\abs{\operatorname{Re}z}<1$ whose real part equals $\abs y$ on the
two boundary lines; the constant $8\Gcat/\pi^2=1-2\nu_4$ is the value
of $h$ at the centre of the strip, and the average of an analytic
function over the circle is its value at the centre. When $B$ and
$\Hilb B$ commute this is the whole argument. In general they do not,
and the new ingredient is a comparison identity for noncommuting
matrices.

We now apply Lemma~\ref{lem:hardy} to the function
\[
U(e^{2\pi\ii t})=e^{-\pi\ii t/2}V(t)^\dagger
\]
built from $V$ and a
partition of the circle. The result is a quadratic form of the type
$L_c$ in \eqref{eq:Lc}, and to write down its coefficients we need two
pieces of notation. First, let
\[
\sigma(x)=\sec\frac{\pi x}2\qquad(0\le x<1),
\]
so that $\sigma(1-x)=\csc(\pi x/2)$ and the coefficients of $\Gd$ in
\eqref{eq:Gdef} are $\beta_k=\sigma(k/d)/(2d^2)$, $1\le k<d$.
Second, for a function $g$ on $[0,1]$ and an
integer $0\le k\le d$, let
\[
\begin{gathered}
P_{d,k}(x)=\binom{d}{k}x^k(1-x)^{d-k},\\
(B_dg)(k)=d\int_0^1P_{d,k}(x)\,g(x)\dd x .
\end{gathered}
\]
The weight $P_{d,k}(x)$ is the probability that a binomial variable
with parameters $(d,x)$ equals $k$; as a function of $x$ it is
concentrated near $x=k/d$ and has integral $1/(d+1)$. Thus
$(B_dg)(k)$ is a weighted average of $g$ over a neighbourhood of the
point $k/d$, normalized up to the factor $d/(d+1)$. We call these Bernstein averages, since the weights are Bernstein
basis polynomials \cite{Bernstein}; $B_d$ is not the classical
Bernstein approximation operator. Finally,
let
\begin{equation}\label{eq:Ddef}
q_k=\frac{(B_d\sigma)(k)}{2d^2},\qquad 1\le k<d,
\end{equation}
which is the coefficient $\beta_k$ of $\Gd$ with the point value
$\sigma(k/d)$ replaced by its Bernstein average, and let $L_q$ be the
form \eqref{eq:Lc} with the coefficients $c_k=q_k$.

\begin{lemma}[Partition forms]\label{lem:partitions}
Let $\ell=(I_0,\dots,I_{d-1})$ be a partition of $[0,1]$ into
consecutive intervals, and let $V(t)=V_j$ for $t\in I_j$. Then
\begin{equation}\label{eq:partitionform}
\begin{split}
\mathcal K_\ell(V)
&:=\norm{P_-[e^{-\pi\ii t/2}V(t)^\dagger]}_2^2-\nu_4\\
&=\frac12\sum_{j<k}\int_{I_j}\!\int_{I_k}
\bigl[\sigma(t-s)(1-\operatorname{Re}z_{jk})\\
&\hspace{4em}+\sigma(1-t+s)\operatorname{Im}z_{jk}\bigr]\dd t\dd s
\;\ge\;0 .
\end{split}
\end{equation}
Extend $V_{j+d}=-\ii V_j$. For a circular partition with $d$ boundary
points, averaging $\mathcal K$ over the $d$ cyclic assignments of
$V_0,\dots,V_{d-1}$ to the arcs gives $L_c$ in \eqref{eq:Lc} with
\begin{equation}\label{eq:arccount}
c_k=\frac1{2d}\int_0^1\sigma(x)\Prob(N_x=k)\dd x,\qquad 1\le k<d,
\end{equation}
where $N_x$ counts boundary points in an arc of length $x$ from a
uniformly random origin. In particular, taking the $d$ boundary points
independent and uniform gives $\Prob(N_x=k)=P_{d,k}(x)$ and hence
$L_q(V)\ge0$.
\end{lemma}

\begin{proof}
The function $U(e^{2\pi\ii t})=e^{-\pi\ii t/2}V(t)^\dagger$ satisfies
$U(z)^4=z^{-1}I$, so Lemma~\ref{lem:hardy} gives
$\mathcal K_\ell\ge0$. For the formula, expand
$\norm{P_-U}_2^2=\langle U,P_-U\rangle$ with the Abel-regularized
kernel of $P_-$, whose limit at argument $x=t-s$ is
$-\tfrac12(1+\ii\cot\pi x)$, and subtract the same expansion at
$V_j=I$, which equals $\nu_4$ by \eqref{eq:nu4}. The diagonal blocks
$I_j\times I_j$ cancel in the difference. A cross block contributes the
kernel times $\tr(V(t)^\dagger V(s))e^{\pi\ii(t-s)/2}-e^{\pi\ii(t-s)/2}$;
combining the block $I_j\times I_k$ with its transpose and using, for
$0<x<1$,
\[
e^{\pi\ii x/2}(1+\ii\cot\pi x)=\frac{\sigma(x)+\ii\,\sigma(1-x)}2,
\]
the two half-squares assemble into the real integrand of
\eqref{eq:partitionform}. The integrals are finite. Near $t-s=1$, the cross-sectional
measure is at most $1-(t-s)$, which cancels the pole of $\sigma(t-s)$.
For adjacent intervals, the region $0<t-s\le x$ is contained in a
triangle of area $x^2/2$; for separated intervals it is empty for small
$x$. Thus its area is $O(x^2)$, and the kernel
$\sigma(1-t+s)=O((t-s)^{-1})$ is integrable. These bounds also justify
the Abel-regularization limit.

For the circular average, first note that in one dimension a single
drop of the phase label by one power of $\ii$ costs nothing: for
$a\in\Z_4$ and a drop at $T$,
\begin{equation}\label{eq:rotatedbranch}
e^{-\pi\ii t/2}\,\ii^{-a+\ind_{\{t>T\}}}
=\ii^{1-a}e^{-\pi\ii T/2}\,U_0([t-T]_1)
\end{equation}
almost everywhere; translation and a constant phase preserve the
moduli of the Fourier coefficients, so the negative-frequency part of
this function has squared norm exactly $\nu_4$.
Consequently $\mathcal K=0$ whenever all $V_j$ are equal multiples of
the identity in the extended sequence, and the circular form, with the value at $V_j=I$
subtracted, is again a combination of the quantities
$1-\operatorname{Re}z_{jk}$ and $\operatorname{Im}z_{jk}$. For a pair
of arcs with labels $j<k$ separated by $k-j$ boundary points along the
direct route, the kernel computation above assigns the coefficient
$\sigma(x)$ to the direct arc and $\sigma(1-x)$ to the complementary
arc. If the
pair straddles the parameter origin, the extension $V_{j+d}=-\ii V_j$
replaces $z_{jk}$ by $-\ii\bar z_{jk}$, which interchanges the two
coefficients---precisely as the direct route then runs through the
complementary arc. Averaging the starting label over the $d$ cyclic
assignments and integrating over a uniform origin yields $L_c$ with
\eqref{eq:arccount}.
Explicitly, put $m=k-j$ and $P_m(x)=\Prob(N_x=m)$. Before simplification,
the coefficient of $1-\operatorname{Re}z_{jk}$ is
\[
\frac1{2d}\int_0^1(1-x)
\bigl[\sigma(x)P_m(x)+\sigma(1-x)P_{d-m}(x)\bigr]\dd x.
\]
Complementary arcs obey $P_{d-m}(x)=P_m(1-x)$. Substituting $x\mapsto1-x$
in the second term combines the weights $1-x$ and $x$, giving
$c_m=(2d)^{-1}\int_0^1\sigma(x)P_m(x)\dd x$.
 For independent uniform boundary points, $N_x$ is
binomial with parameters $(d,x)$, so $\Prob(N_x=k)=P_{d,k}(x)$ and
$c_k=q_k$; nonnegativity of each $\mathcal K$ gives $L_q\ge0$.
\end{proof}

The uniform average $L_q$ does not match $\Gd$: Bernstein averaging
flattens the secant. The next lemma corrects this by giving the
boundary points a non-uniform probability law. Define
\begin{equation}\label{eq:vkernel}
v(x)=-\frac{x(1-x)}2
+\frac{2x-1}{\pi\bigl[\sec\frac{\pi x}2-\csc\frac{\pi x}2\bigr]},
\end{equation}
with removable singularities filled continuously, and
\[
\begin{gathered}
b(x)=-\frac12\Bigl[x(1-x)\sigma(x)-\frac{2x}\pi\Bigr]'',\\
S(x)=\frac{1-z\cot z}{\sin^2z},\qquad
z=\frac\pi2\Bigl(x-\frac12\Bigr),\qquad S(1/2)=\frac13 .
\end{gathered}
\]
The factor $1-x$ cancels the pole of $\sigma(x)$ at $x=1$.
Subtracting $2x/\pi$ removes the regularized endpoint value, leaving a
smooth function that vanishes at both endpoints; in particular $b$ is
bounded on $[0,1]$. Finally put
\begin{equation}\label{eq:es}
\begin{gathered}
e_{d,k}=\sigma(k/d)-(B_d\sigma)(k)-\frac1d(B_db)(k),\\
s_{d,k}=\frac1d(B_dS)(k).
\end{gathered}
\end{equation}

\begin{lemma}[A non-uniform random partition]\label{lem:tilted}
There is a probability density on circular partitions whose average
$\mathcal K_d^v$ of the partition forms $\mathcal K$ is the form
\eqref{eq:Lc} with coefficients
\begin{equation}\label{eq:Kv}
c_k=\frac1{2d^2}\Bigl[(B_d\sigma)(k)
+\frac1d\bigl(B_d(v\sigma)''\bigr)(k)\Bigr],
\end{equation}
and satisfies $\mathcal K_d^v\ge\cstar L_q$ with
$\cstar=\frac32-\frac\pi4$. Moreover,
\begin{equation}\label{eq:matching}
\Gd=\mathcal K_d^v+\frac1{2d^2}L_{e+s},
\qquad s_{d,k}=s_{d,d-k}\ge\frac1{3(d+1)} ,
\end{equation}
where $L_{e+s}$ is the form \eqref{eq:Lc} with coefficients
$c_k=e_{d,k}+s_{d,k}$.
\end{lemma}

\begin{proof}
The function $v$ is symmetric about $1/2$, has zero value and first
derivative at both endpoints, and $v''(0)=v''(1)=(\pi-2)/2$. It has the
expansion
\begin{equation}\label{eq:vFourier}
v(x)=\sum_{n\ge1}a_n\frac{\sin^2(\pi nx)}{\pi^2n^2},
\quad a_n>0,\quad
\sum_{n\ge1}a_n=\frac\pi4-\frac12 .
\end{equation}
To prove this, write $u=x-\tfrac12$, so that
\[
v(x)=\frac{u^2-\tfrac14}2
+\frac u{\sqrt2\,\pi}
\bigl[\csc(\pi u/2)-2\sin(\pi u/2)\bigr].
\]
The nonconstant even Taylor coefficients of $u\csc(\pi u/2)$ are
positive, and for $a=\pi/2$ and $\abs{au}\le\pi/4$,
\[
-[u\sin(au)]^{(4)}=a^3[4\cos(au)-au\sin(au)]>0 ,
\]
so $v^{(4)}>0$ on $(0,1)$. Four integrations by parts then give
\[
a_n=\frac1{4\pi^2n^2}\int_0^1v^{(4)}(x)[1-\cos(2\pi nx)]\dd x>0,
\]
with $a_n=O(n^{-2})$; Fourier inversion and the value $v''(0)$ fix the
sum in \eqref{eq:vFourier}.

Replace the uniform law of the independent boundary points
$T_1,\dots,T_d$ by the law with density, relative to the uniform one,
\begin{equation}\label{eq:partitiondensity}
\begin{split}
\varpi_d&=1+\frac1d\sum_{a<b}v''(T_a-T_b)\\
&=\cstar+\frac1d\sum_{n\ge1}a_n
\Bigl|\sum_{j=1}^de^{2\pi\ii nT_j}\Bigr|^2,
\end{split}
\end{equation}
using the periodic extension of $v$ in the differences; the second form
follows from $v''(x)=2\sum_na_n\cos(2\pi nx)$ and
\eqref{eq:vFourier}. Its expectation is one, and
$\varpi_d\ge\cstar>0$ pointwise. For a single Fourier mode
$v_n=\sin^2(\pi nx)/(\pi^2n^2)$, conditioning on two boundary points
shows that the generating function of the arc count under this law is
\[
(1-x+xz)^d+(d-1)\,v_n(x)\,(z-1)^2(1-x+xz)^{d-2},
\]
and $(z-1)^2(1-x+xz)^{d-2}$ generates exactly
$P_{d,k}''(x)/[d(d-1)]$. Summing the modes with \eqref{eq:vFourier}
gives
\begin{equation}\label{eq:tiltedcount}
\Prob_{\varpi_d}(N_x=k)=P_{d,k}(x)+\frac{v(x)}dP_{d,k}''(x).
\end{equation}
Insert \eqref{eq:tiltedcount} into \eqref{eq:arccount} and integrate by
parts twice: the endpoint terms vanish because $v\sigma$ vanishes at both
endpoints, $(v\sigma)'$ has finite one-sided limits, and $P_{d,k}$
vanishes
there. This gives \eqref{eq:Kv}. The
inequality $\mathcal K_d^v\ge\cstar L_q$ holds because the average under $\varpi_d$
of the pointwise-nonnegative forms $\mathcal K$ is bounded below by $\cstar$
times the uniform average, which is $L_q$.

Finally, differentiating \eqref{eq:vkernel} gives $(v\sigma)''=b-S$,
so the
coefficient of the difference $\Gd-\mathcal K_d^v$ at separation $k$ is
$[e_{d,k}+s_{d,k}]/(2d^2)$, which is \eqref{eq:matching}. The symmetry
$S(x)=S(1-x)$ and $P_{d,k}(x)=P_{d,d-k}(1-x)$ give
$s_{d,k}=s_{d,d-k}$, and the bounds $\sin^2z\le z^2$ and
$1-z\cot z\ge z^2/3$ (a positive Taylor series) give $S\ge1/3$, hence
$s_{d,k}\ge\tfrac13\cdot\tfrac1{d}(B_d1)(k)=\tfrac1{3(d+1)}$, using
$\int_0^1P_{d,k}=1/(d+1)$.
\end{proof}

By \eqref{eq:matching}, the deficit form $\Gd$ dominates
$\cstar L_q$ up to the error term $L_{e+s}/(2d^2)$. The error
coefficients $e_{d,k}$ are small but not symmetric in
$k\leftrightarrow d-k$, so $L_{e}$ is not simply a combination of the
nonnegative quantities $p_{jk}$: by \eqref{eq:Lc},
\[
L_{e}=2\sum_{j<k}e_{d,k-j}\,p_{jk}
+\sum_{j<k}t_{k-j}\operatorname{Im}z_{jk},
\]
where $t_k=e_{d,d-k}-e_{d,k}$, and the skew part multiplies bare
imaginary parts. The next lemma
controls those through triangle ``circulations'', which are themselves
dominated by the forbidden overlaps.

\begin{lemma}[Circulation bound]\label{lem:circulation}
For $i<j<k$, let
$C_{ijk}=\operatorname{Im}z_{ij}+\operatorname{Im}z_{jk}
-\operatorname{Im}z_{ik}$. Then
\begin{equation}\label{eq:circulation}
\abs{C_{ijk}}\le5\,(p_{ij}+p_{jk}+p_{ik}).
\end{equation}
Moreover, with $\omega_{ijk}=t_{j-i}+t_{k-j}-t_{k-i}$,
\begin{equation}\label{eq:Hodge}
L_e=2\sum_{j<k}e_{d,k-j}\,p_{jk}
+\frac1d\sum_{i<j<k}\omega_{ijk}\,C_{ijk}.
\end{equation}
\end{lemma}

\begin{proof}
The sum-of-squares identity \eqref{eq:circulationSOS} in
Appendix~\ref{app:circulation} proves \eqref{eq:circulation}. For
\eqref{eq:Hodge}, compare the coefficients of
$\operatorname{Im}z_{ab}$, $k=b-a$, on both sides. The triangles
containing the edge $(a,b)$ are $(a,b,c)$ with $c>b$, $(c,a,b)$ with
$c<a$, and $(a,c,b)$ with $a<c<b$; they contribute
$\omega_{abc}+\omega_{cab}-\omega_{acb}$ to the coefficient. Each
contains the term $+t_k$, in total $(d-2)t_k$. The remaining terms are
partial sums $\mathsf S(N)=\sum_{m=1}^Nt_m$: they add up to
\[
2t_k+[\mathsf S(d{-}1{-}b)-\mathsf S(d{-}1{-}a)]
+[\mathsf S(a)-\mathsf S(b)],
\]
and the skew symmetry $t_{d-m}=-t_m$ gives
$\mathsf S(d{-}1{-}a)-\mathsf S(d{-}1{-}b)
=\sum_{m=d-b}^{d-1-a}t_m=-\sum_{m=a+1}^bt_m=\mathsf S(a)-\mathsf S(b)$,
so the two brackets cancel. The total coefficient is $d\,t_k$, matching
the left side.
\end{proof}

Two consequences. First, a coarse bound: each edge lies in $d-2$
triangles, so with $b_d=\max_k\abs{e_{d,k}}$ and
$\abs{\omega_{ijk}}\le6b_d$,
\begin{equation}\label{eq:coarseerror}
L_e\ge-2b_d\Fd-\frac{6b_d}d\cdot5\,(d-2)\Fd\ge-32\,b_d\Fd .
\end{equation}
Second, a sharper per-edge bound: define the absolute triangle load at
separation $k$,
\begin{equation}\label{eq:load}
\ell_{d,k}=\sum_{j=1}^{k-1}\abs{t_j+t_{k-j}-t_k}
+\sum_{j=k+1}^{d-1}\abs{t_k+t_{j-k}-t_j},
\end{equation}
which collects $\abs{\omega}$ over the triangles containing one edge of
separation $k$. Then \eqref{eq:circulation} and \eqref{eq:Hodge} give
\begin{equation}\label{eq:fineerror}
L_e\ge2\sum_{j<k}
\Bigl[e_{d,k-j}-\frac5{2d}\,\ell_{d,k-j}\Bigr]p_{jk}.
\end{equation}

It remains to know that the error coefficients $e_{d,k}$---explicit
numbers depending only on $d$ and $k$---are small enough. That is the
content of the following lemma, whose proof---a Stein-type estimate
for the Bernstein averages at large $d$, and an exact interval
computation at small $d$---is given in Appendix~\ref{app:scalar}.

\begin{lemma}[Bounds on the error coefficients]\label{lem:scalar}
(i) For every $d\ge66$,
\begin{equation}\label{eq:largeM}
b_d=\max_{1\le k<d}\abs{e_{d,k}}
\le\frac{45}{32(d+1)(d+2)} .
\end{equation}
(ii) For every $2\le d\le65$ and every $1\le k<d$,
\begin{equation}\label{eq:finitecondition}
\frac1{3(d+1)}+e_{d,k}-\frac5{2d}\,\ell_{d,k}
\ge\frac1{6(d+1)} .
\end{equation}
\end{lemma}

\begin{theorem}[Deficit certificate]\label{thm:certificate}
Let $d\ge2$, $r\ge1$, and let $V_0,\dots,V_{d-1}\in U(r)$ satisfy
$V_j^4=I$. Then $L_q(V)\ge0$ and
\begin{equation}\label{eq:certificate}
\Gd(V)\ge\cstar L_q(V)+\eta_d\Fd(V),
\end{equation}
where
\[
\eta_d=
\begin{cases}
[6d^2(d+1)]^{-1},&2\le d\le65,\\
(2d-131)\,[6d^2(d+1)(d+2)]^{-1},&d\ge66,
\end{cases}
\]
is positive in every dimension.
\end{theorem}

\begin{proof}
Lemma~\ref{lem:partitions} gives $L_q\ge0$. By
Lemma~\ref{lem:tilted},
$\Gd\ge\cstar L_q+L_{e+s}/(2d^2)$, so it suffices to bound
$L_{e+s}\ge2d^2\eta_d\Fd$. The symmetric part is immediate: since
$s_{d,k}=s_{d,d-k}$, definition \eqref{eq:Lc} gives
$L_s=2\sum_{j<k}s_{d,k-j}\,p_{jk}\ge\tfrac2{3(d+1)}\Fd$ by
\eqref{eq:matching}. For $2\le d\le65$, combine this with
\eqref{eq:fineerror} and Lemma~\ref{lem:scalar}(ii):
\[
L_{e+s}\ge2\min_k\Bigl[\frac1{3(d+1)}+e_{d,k}
-\frac5{2d}\ell_{d,k}\Bigr]\Fd
\ge\frac{\Fd}{3(d+1)} .
\]
For $d\ge66$, use instead \eqref{eq:coarseerror} and
Lemma~\ref{lem:scalar}(i):
\[
\begin{split}
L_{e+s}&\ge\Bigl[\frac2{3(d+1)}-\frac{45}{(d+1)(d+2)}\Bigr]\Fd\\
&=\frac{2d-131}{3(d+1)(d+2)}\,\Fd>0 .
\end{split}
\]
Dividing by $2d^2$ gives \eqref{eq:certificate}.
\end{proof}

Combining Lemma~\ref{lem:companion} with
Theorem~\ref{thm:certificate} and the fact that $\Fd \geq 0$ proves the inequality of
Theorem~\ref{thm:main}, and in fact for every local dimension: every
projective strategy on $\Phi_n$, in particular on $\Phi_d$, has
$\IF(d)-I_d=\tfrac{8d}{d-1}\Gd(V)\ge0$.

\section{Proof that the Fourier-phase measurements uniquely saturate the bound}
\label{sec:equality}

\begin{proof}[Proof of Theorem~\ref{thm:main}]
The fact that any strategy obeys $I_d \leq \IF(d)$ was proved in the previous section and is the first part of Theorem~\ref{thm:main}. What remains is to show that if $I_d = \IF(d)$ for a given strategy, then the projectors of that strategy are the Fourier-phase measurements given in \eqref{eq:standard} up to local unitary rotations.

To that end, suppose that $I_d=\IF(d)$. We
keep the local dimension $n$ general through Steps~1--3, which make no
use of $n=d$ and are reused in Appendix~\ref{app:general}; only
Step~4 specializes. Under equality,
$\Gd(V)=0$ by \eqref{eq:normalization}, so $\Fd(V)=0$ by
Theorem~\ref{thm:certificate}, and by \eqref{eq:forbiddenspectral}
every forbidden product vanishes:
\begin{equation}\label{eq:vanishing}
E_{j\ell}E_{k,\ell+u}=0\qquad(j<k,\ \ell\in\Z_4,\ u\in\{1,2\}),
\end{equation}
since a projection product with $\tr(EF)=\tr[(EF)(EF)^\dagger]=0$ is
itself zero.

\emph{Step 1: the tuple commutes.}
Fix $j<k$. By \eqref{eq:vanishing}, $E_{j\ell}E_{km}\ne0$ requires
$m\in\{\ell,\ell-1\}$, and taking adjoints, $E_{km}E_{j\ell'}\ne0$
requires $m\in\{\ell',\ell'-1\}$. For $\ell\ne\ell'$ the two sets
$\{\ell,\ell-1\}$ and $\{\ell',\ell'-1\}$ have at most one element in
common, so at most one value of $m$ can make
$E_{j\ell}E_{km}E_{j\ell'}$ nonzero. Inserting $\sum_mE_{km}=I$,
\[
0=E_{j\ell}E_{j\ell'}=\sum_{m\in\Z_4}E_{j\ell}E_{km}E_{j\ell'}
\qquad(\ell\ne\ell'),
\]
and a sum with at most one nonzero term that equals zero has none.
Hence $E_{j\ell}E_{km}E_{j\ell'}=0$ for all $m$ whenever
$\ell\ne\ell'$. Inserting $\sum_\ell E_{j\ell}=I$ on both sides of
$E_{km}$ therefore gives
\[
E_{km}=\sum_{\ell,\ell'}E_{j\ell}E_{km}E_{j\ell'}
=\sum_\ell E_{j\ell}E_{km}E_{j\ell},
\]
which commutes with every $E_{j\ell}$. Applying this to every pair
$j<k$ makes all the spectral projections $E_{j\ell}$, and hence all
$V_j$, simultaneously diagonalizable.

\emph{Step 2: the tuple is a direct sum of cuts.}
Decompose $\C^{4n}$ into joint eigenspaces, on each of which every
$V_j$ acts as a single phase $\ii^{\ell_j}$. By \eqref{eq:vanishing}, for $j<k$ the
labels satisfy $\ell_k\in\{\ell_j,\ell_j-1\}$: along $j=0,\dots,d-1$
the label is constant except for at most one decrease by one. On each
joint eigenspace, therefore,
\begin{equation}\label{eq:cuts}
V_j=\ii^{\,a-\ind_{\{j>j_0\}}},
\qquad a\in\Z_4,\qquad 0\le j_0<d,
\end{equation}
a configuration we call a \emph{cut}.

\emph{Step 3: every cut gives the standard measurements.}
For $s=d-1-j_0$, the phases $\lambda_j=e^{\pi\ii j/(2d)}$ satisfy
\[
\lambda_j\,\ii^{-\ind_{\{j>j_0\}}}
=\lambda_s^{-1}\lambda_{[j+s]_d},
\]
as one checks separately for $j\le j_0$ and $j>j_0$. Thus the diagonal
$\diag(\lambda_jV_j)$ of a cut equals the standard diagonal
$\diag(\lambda_j)$ with its frequencies shifted by $s$, times the
global phase $\ii^a\lambda_s^{-1}$.
With $D_s=\diag(\omega^{bs})_{b\in\Z_d}$, the shift has the orientation
\begin{equation}\label{eq:cutconjugation}
\mathsf F\diag(\lambda_{[j+s]_d})\mathsf F^\dagger
=D_s^\dagger R_0D_s.
\end{equation}
Each $D_s$ commutes with every $C_b$, and the global phase
$\ii^a\lambda_s^{-1}$ cancels from \eqref{eq:rootmeasurements}.
Consequently, on a joint eigenspace belonging to this cut, the common
intertwiner with the standard PVMs is
$W_s=D_s^\dagger\mathsf F$, by \eqref{eq:fourierconjugation}.
Taking these intertwiners on the joint eigenspaces gives a unitary $W$
on $\C^d\otimes\C^{4n}$ such that
\[
\RootP^t_b(V)=W(\Pi^t_b\otimes I_{4n})W^\dagger
\qquad\text{for all }t,b.
\]

\emph{Step 4: descent to the original space.}
To see that the projectors $\Pi^t_a$ generate $M_d(\C)$, first
conjugate all four PVMs by $\mathsf F$. By
\eqref{eq:fourierconjugation}, their first basis is computational and
their second consists of the columns of
$U_{\rm F}=R_0=\mathsf F\diag(\lambda_j)\mathsf F^\dagger$. An operator
commuting with the first PVM is diagonal, and $U_{\rm F}$ has entries
\[
(U_{\rm F})_{ab}
=\frac{1-\ii}{d\,[1-e^{2\pi\ii(a-b+1/4)/d}]},
\]
all nonzero, so a diagonal operator commuting with the rank-one
projectors onto its columns has all diagonal entries equal. The
commutant of the amplification $\{\Pi^t_a\otimes I_{4n}\}$ is therefore
$I_d\otimes M_{4n}(\C)$: every projection commuting with all four PVMs
is $I_d\otimes p$, its rank is divisible by $d$, and the restriction of
the PVMs to its range is again a standard amplification.

Apply this to the range projection $\Lambda=\Zisom\Zisom^\dagger$ of
the isometry \eqref{eq:reducingisometry}, which commutes with
every $\RootP^t_a(V)$ by \eqref{eq:intertwine}: thus $W^\dagger\Lambda W=I_d\otimes p$ for a
projection $p$ on $\C^{4n}$. Now set $n=d$. Comparing ranks,
$d=\rank\Lambda=d\rank p$ forces $p=vv^\dagger$ for a unit vector
$v\in\C^{4d}$, and $T=(I_d\otimes v^\dagger)W^\dagger\Zisom$ is a
unitary on $\C^d$ with
\[
TP^t_aT^\dagger
=(I_d\otimes v^\dagger)(\Pi^t_a\otimes vv^\dagger)(I_d\otimes v)
=\Pi^t_a ,
\]
which is the classification claimed in Theorem~\ref{thm:main} with
$U=T^\dagger$.

Conversely, if $P^t_a=U\Pi^t_aU^\dagger$, the behaviour coincides with
that of the standard strategy, since
$\tr_d(U\Pi U^\dagger\,U\Pi'U^\dagger)=\tr_d(\Pi\Pi')$, and the
standard strategy attains $\IF(d)$ by Lemma~\ref{lem:companion} at
$V_j=I$.
\end{proof}

\section{Discussion}

For projective measurements on a prescribed maximally entangled state,
maximal CGLMP violation determines the Fourier-phase measurements and
their multiplicity. This settles the first assertion of Problem~27.B
in the projective class. Extending the argument to arbitrary
$d$-outcome POVMs remains a separate problem in general $d$.
Optimization over states is also different: partially entangled states
can give a larger value~\cite{ADGL}.

The certificate proves the quantitative estimate
\begin{equation}\label{eq:overlapstability}
\Fd(V)\le\frac{d-1}{8d\,\eta_d}\bigl(\IF(d)-I_d\bigr).
\end{equation}
This controls forbidden overlaps in the companion representation.
A bound on the distance between measurement families would additionally
require an approximate commuting-cut decomposition, descent through the
reducing isometry, and control of constants as $n$ varies. Neither such
a measurement-distance bound nor a lower bound matching the $O(s/n)$
construction in Appendix~\ref{app:general} is proved here.

Noise resistance must specify both the noise and the witness. For
uniform behaviour noise $q_v=vq+(1-v)q_{\rm unif}$, with
$q_{\rm unif}(a,b\,|\,x,y)=d^{-2}$, one has
\[
I_d(q_{\rm unif})=0,\qquad I_d(q_v)=vI_d(q).
\]
For $I_d(q)>2$, the threshold for violation of this fixed CGLMP
inequality is therefore $v_{\rm CGLMP}=2/I_d(q)$. The maximizing
strategies minimize this threshold. For rank-one PVMs on $\Phi_d$,
white noise on the state induces exactly $q_{\rm unif}$; for unequal
projection ranks the induced noise need not be uniform.
By contrast, the threshold against the entire local polytope
$\mathcal L_d$ is
\[
v_{\mathcal L}(q)=\max\{v\in[0,1]:vq+(1-v)q_{\rm unif}\in\mathcal L_d\}.
\]
Mishra and Senthilkumar \cite[\texttt{clauses/noise-literal},
Theorem~1]{Mishra26b} give better values of $v_{\mathcal L}$ for
$d\ge4$ using projective measurements with unused outcomes. This does
not contradict optimality for the fixed CGLMP witness or establish a
counterexample for rank-one PVMs on $\Phi_d$. Kullback--Leibler
statistical strength, defined by relative entropy to local behaviours
for a specified setting distribution \cite{VDGG}, is another objective;
its maximization does not follow from a CGLMP value bound.

\appendix

\section{Proof of the negative-frequency bound}\label{app:hardy}

This appendix proves Lemma~\ref{lem:hardy}. The scalar function used in
the proof is defined holomorphically on $\abs{\operatorname{Re}z}<1$ by
\[
h(0)=\frac{8G}{\pi^2},\qquad h'(0)=0,\qquad
h''(z)=-\sec\frac{\pi z}{2}.
\]
On the real interval $-1\le x\le1$ it has the absolutely and uniformly
convergent representation
\begin{equation}\label{eq:hseries}
h(x)=\frac8{\pi^2}\sum_{m\ge0}\frac{(-1)^m}{(2m+1)^2}
\cos\frac{(2m+1)\pi x}{2}.
\end{equation}
For real $|x|<1$, the corresponding secant expansion is understood as
the Abel limit
\[
\sec(\pi x/2)=\lim_{r\uparrow1}2\sum_{m\ge0}(-1)^mr^{2m+1}
\cos((2m+1)\pi x/2).
\]
Integrating twice gives
\eqref{eq:hseries}, with $h(\pm1)=0$ and \mbox{$h'(0)=0$}.
Equation~\eqref{eq:hseries} is not a complex-domain series: at any
nonreal $z$ its summands fail to tend to zero. The differential equation
defines the holomorphic extension used below.
For matrices whose spectra lie in the open strip, $h(X)$ is defined by
holomorphic functional calculus, not substitution into a Taylor series
about zero. In particular,
$\tr_r h(X)=r^{-1}\sum_{\lambda\in\spec X}h(\lambda)$, with algebraic
multiplicities; boundary values below are limits of the real part of
this trace.

The mechanism is as follows. Lifting $U$ to
$\mathcal W(s)=sU(s^4)$ produces two commuting Hermitian involutions
$A=\operatorname{Re}\mathcal W+\operatorname{Im}\mathcal W$ and
$B=\operatorname{Re}\mathcal W-\operatorname{Im}\mathcal W$.
Parseval's identity expresses $1-2\norm{P_-U}_2^2$ as
$\int\tr_r(A\Hilb B)$, where $\Hilb$ is the circular Hilbert transform.
Pinching $\Hilb B$ along $B$ preserves this trace pairing. The positive
pinching identity below compares the resulting commuting matrices with
$B+\ii\Hilb B$, whose holomorphic trace can be averaged using its
analytic extension to the disk. The boundary data
$\operatorname{Re}h(\pm1+\ii y)=|y|$ and the mean-value property then
give the bound $h(0)=8G/\pi^2$. The detailed proof includes both the
strip-boundary and disk-boundary limiting arguments.

\emph{Details.---}Throughout, $B$ is a Hermitian matrix with $B^2=I$
and $Y$ is Hermitian. The \emph{pinching} of $Y$ along $B$,
\[
\Pi_B(Y)=\frac{Y+BYB}2,
\]
keeps only the part of $Y$ commuting with $B$; the matrix $C$ above is
$\Pi_B(\Hilb B)$. The proof compares the
matrix $B+\ii Y$ to $B+\ii\,\Pi_B(Y)$ through a nonnegative measure
built from spectra of a non-Hermitian matrix pencil. The first lemma
supplies the moments of that measure; it is
Hein\"avaara's identity \cite[Eq.~(16), proof of Lemma~2.5]{Heinavaara26},
with $A-aI=H^{-1}$, his second matrix equal to $Y$, and the trace divided
by $r$. The holomorphic pinching application is developed below. We write $X_+$ for the
positive part of a Hermitian matrix $X$.

\begin{lemma}[Trace-density moments]\label{lem:density}
Let $H=H^\dagger$ be invertible and $Y=Y^\dagger$ in $M_r(\C)$.
Counting algebraic multiplicities, define
\[
w_{H,Y}(y)=\frac1{2\pi r}
\sum_{\lambda\in\spec(H^{-1}(Y-yI))}\abs{\operatorname{Im}\lambda}.
\]
This function is continuous, nonnegative, and supported on the spectral
interval of $Y$. For every integer $m\ge0$,
\begin{equation}\label{eq:densitymoments}
\begin{split}
\int y^m w_{H,Y}(y)\dd y
&=\frac1{(m+1)(m+2)!}\\
&\quad\times\left.\frac{\mathrm d^{m+2}}{\mathrm dt^{m+2}}
\tr_r(H+tY)_+^{m+1}\right|_{t=0}.
\end{split}
\end{equation}
\end{lemma}

\begin{proof}
For $\operatorname{Im}z>0$, an eigenvector $v$ of $L(z)=H^{-1}(Y-zI)$
with eigenvalue $\lambda$ satisfies $(Y-zI)v=\lambda Hv$; taking
imaginary parts of $v^\dagger(Y-zI)v=\lambda\,v^\dagger Hv$ gives
\[
-\operatorname{Im}z\,\norm v^2
=\operatorname{Im}\lambda\;v^\dagger Hv .
\]
Since $H$ is invertible, the spectrum of $L(z)$ therefore avoids the
real axis, and it splits into an upper and a lower cluster. The sum
\[
g(z)=r^{-1}\sum_{\operatorname{Im}\lambda(L(z))>0}\lambda(L(z))
\]
is holomorphic in the upper half-plane, being the trace of $L(z)$ times
its upper-cluster Riesz projection, with algebraic multiplicities also
at repeated or defective eigenvalues.

The sum extends continuously to the real axis. Choose disjoint small
disks around the limiting eigenvalues. In a connected upper
half-neighbourhood, the number of selected roots in each disk is
constant. A real limiting cluster contributes an integer multiple of
its real eigenvalue. At real $y$, the characteristic polynomial
$\det(Y-yI-\lambda H)/\det H$ has real coefficients, so nonreal roots
occur in conjugate pairs; half of each pair is selected, whence
\[
\operatorname{Im}g(y+\ii0)=\pi\,w_{H,Y}(y).
\]
In particular $w_{H,Y}$ is continuous and nonnegative. Outside the
spectral interval of $Y$, the matrix $Y-yI$ is definite, so the pencil
is similar, up to sign, to the Hermitian matrix
$(\pm(Y-yI))^{1/2}H^{-1}(\pm(Y-yI))^{1/2}$, and the density vanishes.

Near $w=0$, put
\[
K(w)=H^{-1}(I-wY),\qquad q(w)=\tr_r\bigl(K(w)\Pi_-(w)\bigr),
\]
where $\Pi_-(w)$ continues the negative spectral cluster of $H^{-1}$.
For small real $w$, similarity to $(I-wY)^{1/2}H^{-1}(I-wY)^{1/2}$
shows that $q$ has real Taylor coefficients. In the upper exterior
domain, $L(z)=-zK(1/z)$, and the selected upper cluster corresponds to
the continued negative cluster, so $g(z)=-zq(1/z)$. Writing
$q(w)=\sum_jq_jw^j$, Schwarz reflection and Cauchy's formula around the
spectral slit give
\begin{align}\label{eq:Cauchydensity}
&g(z)=-q_0z-q_1+\int\frac{w_{H,Y}(y)}{y-z}\dd y, \nonumber \\
& \int y^mw_{H,Y}(y)\dd y=[w^{m+2}]\,q(w).
\end{align}
The endpoint arcs vanish because pencil eigenvalues remain bounded near
finite points; continuity controls the two banks of the slit.

It remains to identify the Taylor coefficients of $q$ with the
derivatives in \eqref{eq:densitymoments}. Let $k=m+2$ and
$\mathsf R_H(\lambda)=(\lambda I-H)^{-1}$. Contours enclose the
indicated positive or negative spectra and exclude zero. Expanding the
logarithmic determinant ratio that generates the spectral cluster of
$K(w)$ gives
\begin{align*}
[w^k]\,q(w)
&=\frac{(-1)^k}{kr\,2\pi\ii}
\oint_{\Gamma_{H^{-1},-}}
\Tr[((\zeta H-I)^{-1}Y)^k]\dd\zeta\\
&=-\frac1{kr\,2\pi\ii}
\oint_{\Gamma_{H,-}}
\lambda^{k-2}\Tr[(\mathsf R_H(\lambda)Y)^k]\dd\lambda\\
&=\frac1{kr\,2\pi\ii}
\oint_{\Gamma_{H,+}}
\lambda^{k-2}\Tr[(\mathsf R_H(\lambda)Y)^k]\dd\lambda .
\end{align*}
The second line substitutes $\lambda=1/\zeta$; the third uses the decay
$O(\lambda^{-2})$ of the integrand at infinity to trade the negative
contour for the positive one. Positive-cluster functional calculus and
integration by parts, with
\[
\frac{\mathrm d}{\mathrm d\lambda}\Tr[(\mathsf R_H(\lambda)Y)^k]
=-k\Tr[(\mathsf R_H(\lambda)Y)^k\mathsf R_H(\lambda)],
\]
give $[t^k]\,\tr_r(H+tY)_+^{k-1}=(k-1)\,[w^k]\,q(w)$. Together with
\eqref{eq:Cauchydensity}, this proves \eqref{eq:densitymoments}.
\end{proof}

\begin{lemma}[Positive pinching identity]\label{lem:pinching}
Let $B=B^\dagger$ with $B^2=I$, let $Y=Y^\dagger$, and put
$C=\Pi_B(Y)$, $D=Y-C$. There is a nonnegative finite measure $\nu$,
supported on $[-1,1]$ times the spectral interval of $Y$, such that
\begin{equation}\label{eq:pinching}
\begin{split}
\tr_rf(sB+tY)&-\tr_rf(sB+tC)\\
&=t^2\!\int f''(-sx+ty)\dd\nu(x,y)
\end{split}
\end{equation}
for all polynomials $f$. Its density on $\abs x<1$ is
\begin{equation}\label{eq:nudensity}
\frac{\omega(x,y)}{1-x^2},\quad
\omega(x,y)=\frac1{2\pi r}
\sum_{\lambda\in\spec((B-xI)(Y-yI))}\abs{\operatorname{Im}\lambda},
\end{equation}
and
\begin{equation}\label{eq:numajorant}
0\le\frac{\omega(x,y)}{1-x^2}
\le\frac{\norm D}{2\pi\sqrt{1-x^2}},\qquad
\nu(\R^2)=\frac12\tr_r(D^2).
\end{equation}
\end{lemma}

\begin{proof}
Apply Lemma~\ref{lem:density} with $H=(B-xI)^{-1}$, which is Hermitian
and invertible for $\abs x<1$, with the two eigenvalues
$a=(1-x)^{-1}>0$ and $b=-(1+x)^{-1}<0$; note
$H^{-1}(Y-yI)=(B-xI)(Y-yI)$, the pencil in \eqref{eq:nudensity}. Work
in a basis diagonalizing $B$. For an index cycle
$\mathbf i=(i_1,\dots,i_m)$, let $p$ of its entries lie in the $+1$
eigenspace of $B$ and $q=m-p$ in the $-1$ eigenspace, and put
$Y_{\mathbf i}=Y_{i_1i_2}\cdots Y_{i_mi_1}$. The contour form of
\eqref{eq:densitymoments} evaluates by residues at $\lambda=a$:
\begin{equation}\label{eq:twonoderesidue}
\Res_{\lambda=a}
\frac{\lambda^{m-2}}{(\lambda-a)^p(\lambda-b)^q}
=\binom{m-2}{p-1}
\frac{a^{q-1}(-b)^{p-1}}{(a-b)^{m-1}}
\end{equation}
for $p,q\ge1$, $p+q=m$. Substituting the values of $a$ and $b$ gives
the moments of the slice density:
\begin{equation}\label{eq:twonodemoments}
\begin{split}
\int y^{m-2}&\frac{\omega(x,y)}{1-x^2}\dd y\\
&=\frac1{rm\,2^{m-1}}
\sum_{\mathbf i:\,1\le p<m}
\binom{m-2}{p-1}\\
&\quad\times(1-x)^{p-1}(1+x)^{q-1}Y_{\mathbf i}.
\end{split}
\end{equation}
Only mixed cycles, $1\le p<m$, contribute, because pure cycles have
$p=0$ or $q=0$ and no residue at the opposite node.

Now expand the left side of \eqref{eq:pinching} for $f=\exp$. The
coefficient of $t^m$ in the mixed-cycle part of $\tr_re^{sB+tY}$ is, by
the Duhamel expansion,
\[
\frac1{rm}\sum_{\mathbf i:\,1\le p<m}Y_{\mathbf i}
\int_{\substack{u_j\ge0\\\sum_ju_j=1}}
e^{s\sum_jb_{i_j}u_j}\dd u,
\]
where $b_{i}=\pm1$ is the $B$-eigenvalue at index $i$. Pinching by $B$
removes exactly the mixed cycles, so the left side of
\eqref{eq:pinching} is the sum of these terms. Grouping the
positive-$B$ coordinates, whose total weight we call $v$, turns the
simplex integral into a beta integral with density
$v^{p-1}(1-v)^{q-1}/((p-1)!\,(q-1)!)$ against $e^{sv-s(1-v)}$. The
substitution $x=1-2v$, followed by \eqref{eq:twonodemoments}, gives the
coefficient of $t^m$ on the right of \eqref{eq:pinching}, including the
sign in $-sx$. Comparing coefficients of $s$ and $t$ proves the
identity for $f=\exp$ and hence for all polynomials.

For the majorant in \eqref{eq:numajorant}, conjugate the pencil
$(B-xI)(Y-yI)$ by the block-diagonal matrix
\[
\diag(((1+x)/(1-x))^{1/4}I,((1-x)/(1+x))^{1/4}I)
\]
in the eigenbasis of
$B$. The result has Hermitian diagonal blocks and off-diagonal blocks
$\sqrt{1-x^2}\,D_{+-}$ and $-\sqrt{1-x^2}\,D_{+-}^\dagger$, so its
anti-Hermitian part has norm at most $\sqrt{1-x^2}\,\norm D$. By
Bendixson's theorem the same bounds the imaginary part of each
eigenvalue, proving the majorant. Integrating the $m=2$ case of
\eqref{eq:twonodemoments} over $x\in[-1,1]$ gives the total mass
$\tr_r(D^2)/2$.
\end{proof}

\begin{corollary}[Holomorphic pinching on a strip]\label{cor:holomorphic}
Under the hypotheses of Lemma~\ref{lem:pinching}, identity
\eqref{eq:pinching} holds with $s=\rho\in[0,1)$ and $t=\ii$ for every
$f$ holomorphic on $\abs{\operatorname{Re}z}<1$.
\end{corollary}

\begin{proof}
Choose $\rho<\rho'<1$ and $L>\norm Y$. The rectangle
$\{z:\abs{\operatorname{Re}z}\le\rho',\ \abs{\operatorname{Im}z}\le L\}$
contains the spectra of $\rho B+\ii Y$ and $\rho B+\ii C$ and all
points $-\rho x+\ii y$ with $(x,y)$ in the support of $\nu$. Sum the
polynomial identity for a large-parameter resolvent, continue in the
connected exterior of this rectangle, and apply Cauchy's formula on its
boundary.
\end{proof}

\begin{proof}[Proof of Lemma~\ref{lem:hardy}]
\emph{Step 1: a function on the strip.}
On \mbox{$\abs{\operatorname{Re}z}<1$}, define $h$ by
\[
h(0)=\frac{8\Gcat}{\pi^2},\qquad h'(0)=0,\qquad
h''(z)=-\sec\frac{\pi z}2 ;
\]
explicitly $h'(z)=-\frac2\pi\log(\sec\frac{\pi z}2+\tan\frac{\pi z}2)$,
where the logarithm is the unique holomorphic branch on the strip
whose value at $z=0$ is zero.
Its real part $u(x,y)=\operatorname{Re}h(x+\ii y)$ extends continuously
to the closed strip and satisfies
\begin{equation}\label{eq:stripscalar}
u(\pm1,y)=\abs y,\qquad
0\le u(x,y)\le\frac{8\Gcat}{\pi^2}+\abs y .
\end{equation}
Indeed, on the boundary line $x=1$ one computes
\[
-\operatorname{Im}h'(1+\ii y)=\operatorname{sgn}y,
\]
so
$u(1,\cdot)$ has the slopes of $\abs y$, and Taylor's formula with the
classical integral $\int_0^1(1-x)\sec(\pi x/2)\dd x=8\Gcat/\pi^2$ gives
$u(1,0)=0$. Since $h$ is even, $u$ is even in $x$, and
$u_{xx}=-\operatorname{Re}\sec(\pi(x+\ii y)/2)\le0$: $u$ is concave in
$x$, hence at least its boundary values $\abs y\ge0$. On the centre
line, $u(0,y)=8\Gcat/\pi^2+\int_0^y(y-s)\sech(\pi s/2)\dd s
\le8\Gcat/\pi^2+\abs y$ because $\int_0^\infty\sech(\pi s/2)\dd s=1$;
concavity extends the upper bound to the strip.

For a matrix $X$ with spectrum in the closed strip, put
\[
\mathscr U(X)=r^{-1}\sum_{\lambda\in\spec X}\operatorname{Re}h(\lambda).
\]
Apply Corollary~\ref{cor:holomorphic} with $f=h$, take real parts, and
let $\rho\uparrow1$. Since $C=\Pi_B(Y)$ commutes with $B$, the matrix
$B+\ii C$ is normal with eigenvalues $\pm1+\ii c$, where $c$ runs over
the spectrum of $C$ on the two eigenspaces; by the boundary values in
\eqref{eq:stripscalar}, $\mathscr U(B+\ii C)=\tr_r\abs{\Pi_B(Y)}$. The
limit therefore reads
\begin{equation}\label{eq:stripmatrix}
\begin{split}
\mathscr U(B+\ii Y)&-\tr_r\abs{\Pi_B(Y)}\\
&=\int\operatorname{Re}\sec\frac{\pi(x+\ii y)}2\dd\nu(x,y)\ge0,
\end{split}
\end{equation}
the integrand being nonnegative on the strip because $\cos(\pi z/2)$
maps it to the right half-plane. The passage $\rho\uparrow1$ is
dominated: each $x$-slice of the integral is bounded through
\eqref{eq:numajorant} by
$\norm D/(\pi\sqrt{1-x^2})$, using
$\int_\R\operatorname{Re}\sec(\pi(\rho x+\ii y)/2)\dd y=2$; on the left
side, restriction to smaller strips followed by passage to the
endpoints proves convergence for the possibly nonnormal matrix
$B+\ii Y$.

\emph{Step 2: two commuting involutions.}
Set $\mathcal W(s)=sU(s^4)$ on the circle, so $\mathcal W$ is unitary
with $\mathcal W^4=I$, and
\[
A=\operatorname{Re}\mathcal W+\operatorname{Im}\mathcal W,\qquad
B=\operatorname{Re}\mathcal W-\operatorname{Im}\mathcal W .
\]
Because the spectrum of $\mathcal W$ lies in $\{\pm1,\pm\ii\}$, the
commuting Hermitian matrices $\operatorname{Re}\mathcal W$ and
$\operatorname{Im}\mathcal W$ satisfy
\[
(\operatorname{Re}\mathcal W)(\operatorname{Im}\mathcal W)=0
\quad\text{and}\quad
(\operatorname{Re}\mathcal W)^2+(\operatorname{Im}\mathcal W)^2=I,
\]
so
$A$ and $B$ are commuting Hermitian involutions at every point of the
circle. The identity $\mathcal W(\ii s)=\ii\,\mathcal W(s)$ shows that
$A$ and $B$ have mean zero. If $U=\sum_kc_kz^k$, then
$\mathcal W=\sum_kc_ks^{4k+1}$: the frequencies of $\mathcal W$ are
$4k+1$, with the same positive/negative split as those of $U$, and
$\norm{P_-\mathcal W}_2=\norm{P_-U}_2$.

\emph{Step 3: Parseval and the conclusion.}
Let $\Hilb$ be the circular Hilbert transform, with multiplier
$-\ii\operatorname{sgn}(k)$ \cite{Katznelson}. Expanding
$\mathcal W=\tfrac12[(1+\ii)A+(1-\ii)B]$ in Fourier coefficients and
using the frequency split above, Parseval's identity gives
\begin{equation}\label{eq:HardyParseval}
\begin{split}
1-2\norm{P_-U}_2^2
&=\int_{\T}\tr_r(A\,\Hilb B)\dd m\\
&\le\int_{\T}\tr_r\abs{\Pi_B(\Hilb B)}\dd m .
\end{split}
\end{equation}
The inequality holds pointwise: pinching is self-adjoint for the trace
pairing and fixes $A$, since $BAB=AB^2=A$; hence
$\tr(A\,\Hilb B)=\tr\bigl(A\,\Pi_B(\Hilb B)\bigr)
\le\tr\abs{\Pi_B(\Hilb B)}$, because $\norm A\le1$.

The analytic matrix function with boundary values $B+\ii\Hilb B$
vanishes at the disk centre, since $B$ has mean zero. The real part of the holomorphic trace of $h$
of this analytic matrix family is harmonic, so the mean-value property
gives
\[
\int_\T\mathscr U(B+\ii\Hilb B)\dd m=\mathscr U(0)=\frac{8\Gcat}{\pi^2}.
\]
For the boundary passage, use radial limits together with the bound
$0\le\mathscr U(X)\le8\Gcat/\pi^2+\tr_r\abs{\operatorname{Im}X}$, which
follows from \eqref{eq:stripscalar}, Schur triangularization, and
convexity of the absolute value on the diagonal; the radial Hilbert
transforms are uniformly $L^2$-bounded, hence the displayed quantities
are uniformly integrable. Applying \eqref{eq:stripmatrix} pointwise in
\eqref{eq:HardyParseval} now gives
$1-2\norm{P_-U}_2^2\le8\Gcat/\pi^2$, which is \eqref{eq:Hardy}.
\end{proof}

\section{The circulation sum of squares}\label{app:circulation}

This appendix proves \eqref{eq:circulation}. Fix $i<j<k$, and write
\[
\begin{array}{lll}
u_\ell=E_{i\ell}E_{j,\ell+2},&
v_\ell=E_{j\ell}E_{k,\ell+2},&
z_\ell=E_{i\ell}E_{k,\ell+2},\\[2pt]
a_\ell=E_{i\ell}E_{j,\ell+1},&
b_\ell=E_{j\ell}E_{k,\ell+1},&
c_\ell=E_{i\ell}E_{k,\ell+1},
\end{array}
\]
the forbidden products on the three edges of the triangle. Let $P$ be
the sum of the squared norms $\norm\cdot_2^2$ of all $a_\ell,b_\ell,
c_\ell$, and $O$ the corresponding sum for $u_\ell,v_\ell,z_\ell$, so
that $P+O=p_{ij}+p_{jk}+p_{ik}$ by \eqref{eq:forbiddenspectral}.

\begin{table}[!htbp]
\centering
\caption{Terms in the circulation identity. Spectral indices are read
modulo four.}
\label{tab:circulation}
\renewcommand{\arraystretch}{1.35}
\begin{tabular}{@{}ccc@{}}
\toprule
$X_{\ell\nu}$ & $Y_{\ell\nu}$ & $\epsilon_\nu$\\
\midrule
$z_\ell$ & $a_\ell$ & $1$\\
$z_\ell^\dagger$ & $b_{\ell+1}^\dagger$ & $1$\\
$v_{\ell+1}$ & $a_\ell^\dagger$ & $1$\\
$u_\ell$ & $c_\ell$ & $-1$\\
$u_\ell^\dagger$ & $b_{\ell+2}$ & $1$\\
$v_{\ell+3}^\dagger$ & $c_\ell^\dagger$ & $-1$\\
\bottomrule
\end{tabular}
\end{table}

Projection completeness ($\sum_mE_{jm}=I$, inserted between the two
factors of $z_\ell$ and $c_\ell$) gives
\begin{equation}\label{eq:Cexpansion}
\begin{split}
C_{ijk}={}&{-2}\sum_\ell\norm{z_\ell}_2^2
+2\operatorname{Re}\sum_\ell
\bigl[\ip{u_\ell}{z_\ell}
+\ip{v_\ell^\dagger}{z_\ell^\dagger}\bigr]\\
&+2\operatorname{Re}\sum_{\ell,\nu}
\epsilon_\nu\ip{X_{\ell\nu}}{Y_{\ell\nu}},
\end{split}
\end{equation}
with the six triples $(X_{\ell\nu},Y_{\ell\nu},\epsilon_\nu)$ of
Table~\ref{tab:circulation}. For $\sigma=\pm1$, completing squares
yields the identity
\begin{equation}\label{eq:circulationSOS}
\begin{split}
2P+5O-\sigma C_{ijk}
={}&\sum_{\ell,\nu}
\norm{Y_{\ell\nu}-\sigma\epsilon_\nu X_{\ell\nu}}_2^2\\
&+3\sum_\ell\norm{u_\ell-\tfrac\sigma3z_\ell}_2^2\\
&+3\sum_\ell\norm{v_\ell^\dagger-\tfrac\sigma3z_\ell^\dagger}_2^2\\
&+\Bigl(\frac73+2\sigma\Bigr)
\sum_\ell\norm{z_\ell}_2^2\;\ge\;0,
\end{split}
\end{equation}
whence $\abs{C_{ijk}}\le2P+5O\le5(P+O)$, which is
\eqref{eq:circulation}. Both \eqref{eq:Cexpansion} and
\eqref{eq:circulationSOS} are finite identities in the twelve
projections: they are verified by expanding
$E_{j3}=I-E_{j0}-E_{j1}-E_{j2}$ at each site, reducing adjacent
same-site products by orthogonality, and comparing coefficients of
cyclic words---choosing the lexicographically least cyclic rotation of
each word and splitting each real part into the half-sum of a word and
its adjoint.



\section{Bounds on the error coefficients}\label{app:scalar}

This appendix proves Lemma~\ref{lem:scalar}. Throughout, $y=1-k/d$,
and
\[
r_{\rm reg}(y)=\csc\frac{\pi y}2-\frac2{\pi y},\qquad
\varphi(y)=y(1-y)\,r_{\rm reg}(y).
\]
The cosecant expansion gives
\begin{equation}\label{eq:rn}
\begin{gathered}
r_{\rm reg}(y)=\sum_{n\ge1}r_ny^{2n-1},\\
r_n=\frac{2(2^{2n-1}-1)\abs{B_{2n}}}{(2n)!}
\Bigl(\frac\pi2\Bigr)^{2n-1},
\end{gathered}
\end{equation}
with Bernoulli numbers $B_{2n}$ and $0<r_n<3\cdot4^{-n}$. Since
$\sigma(x)=\csc(\pi y/2)$ and $1-y=x$, one checks
$x(1-x)\sigma(x)-2x/\pi=\varphi(y)$, so the kernel $b$ of
\eqref{eq:es} is
$b=-\varphi''(y)/2$. Here $x(1-x)$ already cancels the secant pole;
subtracting $2x/\pi$ leaves the contribution of the regular part
$r_{\rm reg}$, with zero endpoint values. The
Bernstein averages themselves are beta expectations: with
$X\sim\Beta(k+1,d-k+1)$ and $Y=1-X$,
$(B_dg)(k)=\tfrac d{d+1}\Ex[g(X)]$, the moments are
\[
\frac d{d+1}\Ex[Y^j]=m_j:=\frac d{d+1}
\frac{(d-k+1)_j}{(d+2)_j},
\]
and $\Ex[Y^{-1}]=(d+1)/(d-k)$, which cancels the pole of
$\sigma(k/d)$ in
\eqref{eq:es}. Substituting \eqref{eq:rn} term by term gives the
absolutely convergent series
\begin{equation}\label{eq:finiteeseries}
\begin{split}
e_{d,k}=\sum_{n\ge1}r_n
\Bigl[&y^j-\Bigl(1+\frac{(j+1)(j+2)}{2d}\Bigr)m_j\\
&+\frac{j(j+1)}{2d}m_{j-1}\Bigr],
\end{split}
\end{equation}
with $j=2n-1$. Each bracket is bounded in absolute value by
$2+4n^2/(d+1)$, since $y^j,m_j\in[0,1]$ and $m_j\le d/(d+1)$.

\subsection{The beta--Stein estimate: proof of (i)}

Let $Z\sim\Beta(dy,d(1-y))=\Beta(d-k,k)$, so $\Ex Z=y$. Multiplying the
$\Beta(d-k,k)$ density by $z(1-z)$ reproduces the $\Beta(d-k+1,k+1)$
kernel, which turns the Bernstein average into a size-biased
expectation,
$(B_dg)(k)=\mbox{$\Ex[Z(1-Z)g(1-Z)]/[y(1-y)]$}$, and integration by parts against
the beta density gives the Stein identity
\cite{GoldsteinReinert,Dobler}
\[
\Ex[(Z-y)g(Z)]=\frac1d\,\Ex[Z(1-Z)g'(Z)].
\]
Combining the two with \eqref{eq:es} yields the exact representation
\begin{equation}\label{eq:betaerror}
e_{d,k}=\frac{\varphi(y)-\Ex\varphi(Z)
+\tfrac12\Ex[(Z-y)\varphi'(Z)]}{y(1-y)} .
\end{equation}
The numerator is a trapezoidal remainder: since
\[
\Ex[(Z-y)\varphi'(y)]=0
\]
may be added freely, with $\xi=Z-y$,
\[
\varphi(y)-\varphi(y+\xi)
+\frac\xi2[\varphi'(y+\xi)+\varphi'(y)]
=\frac{\xi^3}{12}\varphi'''(y)+\epsilon_\xi,
\]
where
$\abs{\epsilon_\xi}\le\norm{\varphi^{(4)}}_\infty\abs\xi^4/24$.
The derivatives of $\varphi$ are controlled by the positive series
\eqref{eq:rn}: for $j=2n-1$, the third and fourth derivatives of
$y^{j+1}-y^{j+2}$ are bounded on $[0,1]$ by $3j(j+1)$ and
$4j(j-1)(j+1)$, and the sums
\[
r_{\rm reg}''(1)+2r_{\rm reg}'(1)=\frac{\pi^2}4,\qquad
r_{\rm reg}'''(1)+3r_{\rm reg}''(1)=\frac{3\pi^2}4
\]
give $\norm{\varphi'''}_\infty\le3\pi^2/4$ and
$\norm{\varphi^{(4)}}_\infty\le3\pi^2$. The central beta moments are
\[
\Ex\xi^3=\frac{2y(1-y)(1-2y)}{(d+1)(d+2)},
\]
\[
\Ex\xi^4=\frac{3y(1-y)[2(1-2y)^2+y(1-y)(d+2)]}{(d+1)(d+2)(d+3)}.
\]
Inserting these into \eqref{eq:betaerror}, with $q=\abs{1-2y}$,
\[
\abs{e_{d,k}}\le\frac{\pi^2}{8(d+1)(d+2)}
\Bigl[q+\frac{3(d+2)}{4(d+3)}
-\frac{3(d-6)}{4(d+3)}q^2\Bigr].
\]
For $d\ge66$, maximizing the bracket over $q\in[0,1]$ gives at most
\[
\frac{3(d+2)}{4(d+3)}+\frac{d+3}{3(d-6)}<\frac98,
\]
the difference from $9/8$ being
\[
(d^2-57d-342)/[24(d-6)(d+3)]>0.
\]
Using $\pi^2<10$ gives
\eqref{eq:largeM}.

\subsection{The interval computation: proof of (ii)}

For $2\le d\le65$, the inequality \eqref{eq:finitecondition} is a
finite set of rational comparisons, verified in exact interval
arithmetic \cite{Moore} as follows. Truncating
\eqref{eq:finiteeseries} after $32$ terms has error at most
\begin{equation}\label{eq:finitetail}
T_d=3\sum_{n=33}^\infty4^{-n}\Bigl(2+\frac{4n^2}{d+1}\Bigr),
\end{equation}
evaluated exactly with
\[
\begin{gathered}
\sum_{n\ge J}q^n=\frac{q^J}{1-q},\\
\sum_{n\ge J}n^2q^n
=q^J\Bigl[\frac{J^2}{1-q}+\frac{2Jq}{(1-q)^2}
+\frac{q(1+q)}{(1-q)^3}\Bigr]
\end{gathered}
\]
at $q=1/4$. For the coefficients, put
\[
A(a,N)=\sum_{i=0}^{N-1}(-1)^i/[(2i+1)a^{2i+1}];
\]
Machin's identity
gives the rational enclosure
\[
\begin{split}
\pi\in{}&16[A(5,128),A(5,129)]\\
&-4[A(239,128),A(239,129)].
\end{split}
\]
Generate the Bernoulli numbers by $B_0=1$ and
$B_m=-(m+1)^{-1}\sum_{i=0}^{m-1}\binom{m+1}iB_i$ through $m=64$, and
substitute the enclosure of $\pi$ into \eqref{eq:rn}. The moments
follow the recursion $m_0=d/(d+1)$,
$m_j=m_{j-1}(d-k+j)/(d+j+1)$. Endpoint inclusion in these rational
operations, followed by adding $[-T_d,T_d]$, encloses each $e_{d,k}$;
the skew coefficients $t_k=e_{d,d-k}-e_{d,k}$ and the loads
$\ell_{d,k}$ of \eqref{eq:load} are enclosed likewise, bounding each
absolute value $\abs{[l,u]}$ by $\max(\abs l,\abs u)$. The parameters
are collected in Table~\ref{tab:intervals}: over all $2\le d\le65$ and
$1\le k<d$, the normalized margin satisfies
\[
(d+1)\Bigl[\frac1{3(d+1)}+\min_k e_{d,k}
-\frac5{2d}\max_k\ell_{d,k}\Bigr]
\ge\frac{1935}{10^4}>\frac16,
\]
which is \eqref{eq:finitecondition}.
The accompanying \path{verification/verify_coefficients.py} checks all
2,080 pairs using Python~3 and only its standard library. Run
\texttt{python3} \path{verification/verify_coefficients.py} from the manuscript
directory. The worst dimension is $d=7$, and the certified minimum
normalized margin exceeds $0.19354436777798498$. The Machin enclosure is
rounded outward to 80 decimal places and the $r_n$ enclosures to 60;
all rounding, propagation, and comparisons use exact rational numbers.
The companion script \path{verification/verify_circulation.py}
checks \eqref{eq:Cexpansion} and both signs of
\eqref{eq:circulationSOS} by exact cyclic-word arithmetic.
Commands, provenance, expected output, and recorded outputs accompany
the scripts in \texttt{verification/}. These computations certify the
stated finite inequalities and identities, not the full analytic proof.

\begin{table}[!htbp]
\centering
\caption{Data for the finite interval computation. The margin uses
lower endpoints for $e_{d,k}$ and upper bounds for $\ell_{d,k}$.}
\label{tab:intervals}
\begin{tabular}{@{}l@{\hspace{1em}}l@{}}
\toprule
Quantity & Data\\
\midrule
Outcome counts & $2\le d\le65$\\
Separations & $1\le k<d$\\
Cosecant terms & $1\le n\le32$\\
Moment indices & $0\le j\le63$\\
Bernoulli indices & $0\le m\le64$\\
Arctangent partial sums & $128,\ 129$ terms\\
Outward rounding ($\pi$, $r_n$) & 80, 60 decimal places\\
Worst dimension & $d=7$\\
Tail start in \eqref{eq:finitetail} & $J=33$\\
Lower bound for the normalized margin & $1935/10000>1/6$\\
\bottomrule
\end{tabular}
\end{table}

\section{Arbitrary local dimension}\label{app:general}

Problem~27.B concerns the maximally entangled state of two $d$-level
systems, which is Theorem~\ref{thm:main}. The proof, however, works on
$\Phi_n$ for every $n$ until its very last step, and yields the
following extension.

\begin{theorem}[General local dimension]\label{thm:general}
Let $d\ge2$ and $n\ge1$. Every projective strategy on $\Phi_n$
satisfies $I_d\le\IF(d)$. Equality holds if and only if $d$ divides
$n$ and there is a unitary $U$ on $\C^n=\C^d\otimes\C^{n/d}$ such that
the tracial PVMs are $P^t_a=U(\Pi^t_a\otimes I_{n/d})U^\dagger$ for
all $t,a$. In particular, if $d$ does not divide $n$, the value
$\IF(d)$ is not attained on $\Phi_n$.
\end{theorem}

\begin{proof}
The inequality was proved for every $n$ at the end of
Section~\ref{sec:partition}. Suppose equality. Steps~1--3 of the
proof of Theorem~\ref{thm:main} make no use of $n=d$ and give a
unitary $W$ on $\C^d\otimes\C^{4n}$ with
$\RootP^t_a(V)=W(\Pi^t_a\otimes I_{4n})W^\dagger$ for all $t,a$. As
in Step~4 there, the commutant argument gives
$W^\dagger\Lambda W=I_d\otimes p$ for the range projection
$\Lambda=\Zisom\Zisom^\dagger$ of the isometry
\eqref{eq:reducingisometry} and a projection $p$ on $\C^{4n}$.
Comparing ranks, $n=\rank\Lambda=d\rank p$, which forces $d\mid n$.
Writing $k=n/d$ and choosing an isometry $\iota:\C^k\to\C^{4n}$ onto
the range of $p$, the unitary
$T=(I_d\otimes\iota^\dagger)W^\dagger\Zisom:
\C^n\to\C^d\otimes\C^k$ satisfies
\[
TP^t_aT^\dagger
=(I_d\otimes\iota^\dagger)(\Pi^t_a\otimes p)(I_d\otimes\iota)
=\Pi^t_a\otimes I_k,
\]
the claimed classification with $U=T^\dagger$.

Conversely, an amplified strategy
$P^t_a=U(\Pi^t_a\otimes I_k)U^\dagger$ on $\Phi_{dk}$ has the same
behaviour as the standard strategy on $\Phi_d$, since
$\tr_{dk}\bigl((\Pi\otimes I_k)(\Pi'\otimes I_k)\bigr)
=\tr_d(\Pi\Pi')$, and therefore attains $\IF(d)$.
\end{proof}

For fixed $d,n$, the space of four $d$-outcome PVMs on $\C^n$ is
compact: it is a closed, bounded subset of a finite-dimensional matrix
space. The Bell functional is continuous, so its optimum is attained.
The equality classification therefore gives a strictly positive gap
when $d\nmid n$, without specifying a rate. Write $s=n\bmod d>0$. Pad $\lfloor n/d\rfloor$ standard blocks with $s$
one-dimensional blocks on which all four PVMs assign one fixed
outcome achieving the local value $2$. The resulting behaviour is the
mixture of the standard behaviour, with weight $d\lfloor n/d\rfloor/n$,
and a deterministic local behaviour, with weight $s/n$, so that
\[
\sup\{I_d\ \text{on}\ \Phi_n\}\ \ge\ \IF(d)-\frac sn\bigl(\IF(d)-2\bigr).
\]
The deficit of the optimum is therefore at most of order $1/n$ along
every residue class; lower bounds on it require a quantitative version
of the rigidity argument and are not pursued here.

% Reference list: entries set ragged right; identifiers break only after "/".
\makeatletter
\def\UrlBreaks{\do\/}\def\UrlBigBreaks{}%
\def\UrlOrds{\do\*\do\-\do\~\do\'\do\"\do\:\do\.}%
\let\oqp@BibitemOpen\BibitemOpen
\def\BibitemOpen{\oqp@BibitemOpen\rightskip\z@\@plus1fil\relax}
\makeatother
\bibliography{oqp27-note-references}

\end{document}
